%% toc: P15 164 164.12
\chapter{Systematic Power in Hempel and Oppenheim: Data, Content and Range, and the Generalized Theory}
\label{sec:p15-w02c-p15-029-hempel-oppenheim}

Part~IV of Hempel and Oppenheim's 1948 paper builds a measure of
systematic power in two stages: a counting definition on minimal
sentences, and a measure-based definition that turns out to be the dual
of logical probability \cite[pp.~164--173]{HempelOppenheim1948Explanation}.
The definitions and theorems are in \cref{sec:p15-systematic-power}; the
computations in the model languages \(L_1\) and \(L_2\) are in
\cref{sec:p15-w02c-p15-028-model-languages}; the comparison with
relevance measures is in \cref{sec:p15-w02b-p15-025-comparison}. This
chapter presents the parts of the source that those chapters do not
contain: what the measure is meant to explicate and how the data are
construed; the source's reading of the two normal forms as indicators of
content and of range, with the duality of minimal and maximal sentences;
its remarks on the uniform measures \(r_1\) and \(g_1\) and on Carnap's
criticism of \(r_1\); the variants of the generalized measure that it
names and sets aside, with its closing evaluation; and a complete proof
of the converse half of its Dualism Theorem (9.13). Each source report is
followed by an exact statement of what it amounts to in the handbook's
notation.

\paragraph{Inputs and conclusions.}
The input is a fixed language as in \cref{def:p15-power-language}, with
ground atoms \(A_1,\ldots,A_n\), assignments \(v\in\Omega=\{0,1\}^n\),
state descriptions \(D_v\), minimal sentences \(M_v\), truth sets \([U]\)
and content families \(\mathcal C(U)\); the original measure \(s_0\) of
\cref{def:p15-power-optimal-split}; the uniform measures \(r_1,g_1\) of
\cref{def:p15-power-uniform-measures}; range measures \(r\) and the
functions \(p_r,s_r\) of \cref{def:p15-power-range-measure} and
\cref{def:p15-power-dual-measures}; and the dual operation \(D\) of
\cref{def:p15-power-dual-syntax}. The conclusions are
\cref{thm:p15-w02c-p15-029-ratio,thm:p15-w02c-p15-029-renaming,thm:p15-w02c-p15-029-duals,thm:p15-w02c-p15-029-inverse,thm:p15-w02c-p15-029-uniform-independence,thm:p15-w02c-p15-029-extension,thm:p15-w02c-p15-029-strict-probability,thm:p15-w02c-p15-029-dualism-converse},
each proved in the annex.

\section{What the measure explicates, and what counts as a datum}
\label{sec:p15-w02c-p15-029-data}

Hempel and Oppenheim start from an intuitive comparison. Some theories
seem powerful because they let many data be derived from little initial
information; others seem weaker because they need more initial data or
yield fewer results. They ask whether such comparisons of explanatory or
predictive power can be made precise by a general numerical measure
\cite[p.~164]{HempelOppenheim1948Explanation}. Their answer is that the
systematic power of a theory \(T\) should reflect the ratio of the amount
of information derivable by means of \(T\) to the amount of initial
information the derivation requires. Since that ratio depends on the body
of information to which \(T\) is applied, the concept is relativized to a
finite class \(K\) of data \cite[pp.~164--165]{HempelOppenheim1948Explanation}.

\begin{hbconvention}{def:p15-w02c-p15-029-data}{Potential data and the theory argument}
Every singular sentence of the language is said to express a potential
datum, and \(K\) is accordingly a finite class of singular sentences.
The theory \(T\) may be any sentence of the language, whether essentially
generalized or not; Hempel and Oppenheim adopt this liberal convention
for the sake of the generality and simplicity of the definitions and
theorems that follow \cite[p.~165]{HempelOppenheim1948Explanation}.

In their note~33, ``datum'' covers potential as well as actual data. The
convention that every singular sentence expresses a potential datum is
plausible especially when the primitive predicates refer to attributes
whose presence or absence in an instance can be ascertained by direct
observation; each singular sentence then describes a logically possible
state of affairs that observation might establish. The assumption that
the primitives express directly observable attributes is, they add, not
essential to the definition or to the formal theory of systematic power
in their sections~8 and~9 \cite[p.~165, n.~33]{HempelOppenheim1948Explanation}.
\end{hbconvention}

The definition (8.1) does not use the ratio of derived to initial
information directly. To keep values between \(0\) and \(1\), Hempel and
Oppenheim first try the fraction of the sentences of \(K\) derivable from
the rest by means of \(T\), and then apply the same fraction to minimal
sentences \cite[pp.~165--166]{HempelOppenheim1948Explanation}. The next
result shows that this normalization does not change the comparison the
intuitive ratio was meant to make.

\begin{hblemma}{thm:p15-w02c-p15-029-ratio}{proved=annex:p15-w02c-p15-029-ratio}{The bounded fraction and the derived-to-initial ratio order theories alike}
Let \(K\) be a finite set of singular sentences whose conjunction \(S\) is
not valid, let \(K'=\mathcal C(S)\), and for a sentence \(T\) let
\(K_T=\{M\in K':T\entails M\}\) be the optimal second part of
\cref{thm:p15-power-optimal-split}. Define the derived-to-initial ratio
\[
 q(T,K)=\frac{|K_T|}{|K'\setminus K_T|},
 \qquad q(T,K)=+\infty\ \text{when }K_T=K'.
\]
Then \(q(T,K)=s_0(T,K)/(1-s_0(T,K))\) whenever \(s_0(T,K)<1\), and for
any two sentences \(T,T'\)
\[
 s_0(T,K)\le s_0(T',K)\quad\Longleftrightarrow\quad q(T,K)\le q(T',K).
\]
In the source's one-name example
(\cref{thm:p15-power-source-examples}), \(s_0=1/3\) and \(q=1/2\): one
minimal sentence is derived from the two used as premises. For the
regrouped data \(K_1,K_2,K_3\) of the same theorem, \(s_0=175/252\) and
\(q=175/77=25/11\).
\end{hblemma}

The lemma uses the optimal partition, whose first part
\(K'\setminus K_T\) is the class of minimal sentences that serve as
premises. The two quantities rank theories identically on fixed data; the
fraction is bounded and the ratio is not, which is why the bounded form
can be compared with a probability.

\begin{hbtheorem}{thm:p15-w02c-p15-029-renaming}{proved=annex:p15-w02c-p15-029-renaming}{The formal theory does not depend on what the predicates refer to}
Let \(L\) and \(L'\) be languages as in \cref{def:p15-power-language},
and let \(\sigma\) be a bijection from the individual constants of \(L\)
onto those of \(L'\) and from the predicate symbols of \(L\) onto those
of \(L'\) that preserves arity. Extend \(\sigma\) to sentences by
replacing every symbol. Then:
\begin{enumerate}
\item \(U\entails V\) in \(L\) if and only if \(\sigma U\entails\sigma V\)
  in \(L'\); in particular \(\sigma\) preserves validity, and maps the
  ground atoms of \(L\) bijectively onto those of \(L'\);
\item for every sentence \(T\) and every finite set \(K\) of singular
  sentences with nonvalid conjunction,
  \(s_0(\sigma T,\sigma K)=s_0(T,K)\);
\item if \(r\) is a range measure on \(L'\), then \(r\circ\sigma\) is a
  range measure on \(L\), and
  \(s_{r\circ\sigma}(T,S)=s_r(\sigma T,\sigma S)\) and
  \(p_{r\circ\sigma}(T,S)=p_r(\sigma T,\sigma S)\) whenever either side
  is defined.
\end{enumerate}
\end{hbtheorem}

The theorem is the exact content of the last sentence of note~33: no
value of \(s_0\), \(s_r\) or \(p_r\) can change when the predicates are
reread as naming observable or unobservable attributes, because the
values depend only on the logical form of the sentences and on the
measure. Whether the singular sentences of a given language can in fact
be checked by observation is a question about applying the theory; it
belongs with the open problems Hempel and Oppenheim list at the end of
Part~IV (\cref{sec:p15-w02c-p15-029-evaluation}).

\section{Content, range and the duality of the two normal forms}
\label{sec:p15-w02c-p15-029-content-range}

Hempel and Oppenheim take the state description, which they also call a
maximal sentence, as the basic notion for measuring logical probability,
as the minimal sentence is basic for systematic power. They describe a
maximal sentence as the dual of a minimal sentence, referring to Church's
dictionary article on formal logic for the notion of duality
\cite[p.~168, n.~36]{HempelOppenheim1948Explanation}. The lists of
minimal and maximal sentences of \(L_1\) and the source's pairs of normal
forms are in \cref{thm:p15-w02c-p15-028-l1-forms}. The duality is exact
in the following sense. For \(v\in\Omega\) write \(\bar v\) for the
assignment with \(\bar v_i=1-v_i\), and treat the ground atoms as the
metavariables of \cref{def:p15-power-dual-syntax}.

\begin{hbtheorem}{thm:p15-w02c-p15-029-duals}{proved=annex:p15-w02c-p15-029-duals}{Minimal and maximal sentences are duals, and so are the two normal forms}
\begin{enumerate}
\item For every \(v\in\Omega\), \(D(M_v)=D_{\bar v}\) and
  \(D(D_v)=M_{\bar v}\). Since \(v\mapsto\bar v\) is a bijection of
  \(\Omega\), the dual operation maps the \(2^n\) minimal sentences
  onto the \(2^n\) maximal sentences and back.
\item For every singular sentence \(U\) written as a Boolean combination
  of the ground atoms,
  \[
   D(U)\equiv\neg U(\neg A_1,\ldots,\neg A_n),\qquad
   [D(U)]=\{\bar v:v\notin[U]\}.
  \]
\item The dual of the minimal normal form of \(U\) is, term by term, the
  maximal normal form of \(D(U)\):
  \[
   D\Bigl(\bigwedge_{v\notin[U]}M_v\Bigr)=\bigvee_{v\notin[U]}D_{\bar v}
   \equiv D(U),
   \qquad |[D(U)]|=|\mathcal C(U)|.
  \]
\end{enumerate}
In \(L_1\) (\cref{def:p15-w02c-p15-028-languages}) this gives the
following. Since \(D(P(a))=P(a)\), the dual of the minimal normal form
\(M_{00}\land M_{01}\) of \(P(a)\) is \(D_{11}\lor D_{10}\), the maximal
normal form of \(P(a)\) itself. For \(U=\neg P(a)\lor Q(a)\), the
expansion of \(P(a)\to Q(a)\), the minimal normal form is \(M_{10}\), and
its dual is \(D_{01}=\neg P(a)\land Q(a)\). That is the maximal normal
form of \(D(U)=\neg P(a)\land Q(a)\), not of \(U\), whose maximal normal
form has three members. The minimal normal form \(\top\) of a valid
sentence dualizes to \(\bot\), the maximal normal form of a contradictory
sentence.
\end{hbtheorem}

The \(L_1\) contrast shows that duality and the content--range
relation are different maps. Duality sends the minimal normal form of
\(U\) to the maximal normal form of another sentence \(D(U)\). The
relation \(|[U]|+|\mathcal C(U)|=2^n\) of
\cref{thm:p15-power-normal-forms} connects the two normal forms of the
same sentence.

\paragraph{The two normal forms as indicators.}
According to Hempel and Oppenheim, the minimal normal form indicates the
content of a singular sentence, because it writes the sentence as a
conjunction of canonical components, no two of which overlap in content
and each of which has the least content possible. The maximal normal form indicates its range:
intuitively, the variety of its possible realizations, or of the possible
states of the world that would make it true
\cite[pp.~168--169]{HempelOppenheim1948Explanation}. Each maximal
sentence describes one possible state of the world as completely as the
language allows, and the state descriptions in the maximal normal form of
a sentence list exactly the possible states that would make it true. The
ranges of two distinct maximal sentences are mutually exclusive, since no
possible state makes two incompatible sentences true
\cite[p.~169]{HempelOppenheim1948Explanation}. For a fuller discussion of
range, and of range and content, they refer to Carnap's 1945 paper on
inductive logic, section~2, and to his \emph{Introduction to Semantics},
sections~18 and~19 \cite[p.~169, n.~37]{HempelOppenheim1948Explanation}.

Range and content, they continue, move in opposite directions: a sentence
that says more is realized in fewer possible ways. In the
normal forms this appears as the count relation \(l_U=m-m_U\) of
\cref{def:p15-w02c-p15-028-counts}. To keep a content measure between
\(0\) and \(1\) they divide the count \(m_U\) by \(m\), which is (9.1),
and they assign to a finite class \(K\) of singular sentences the content
\(g_1(K)=g_1(S)\), where \(S\) is the conjunction of the members of
\(K\) \cite[p.~169]{HempelOppenheim1948Explanation}. The following
theorem states the inverse variation as an order relation and records
what a state description does and does not settle.

\begin{hbtheorem}{thm:p15-w02c-p15-029-inverse}{proved=annex:p15-w02c-p15-029-inverse}{Content and range vary inversely, and what a state description settles}
For singular sentences \(U,V\):
\begin{enumerate}
\item \(U\entails V\) if and only if \([U]\subseteq[V]\), if and only if
  \(\mathcal C(V)\subseteq\mathcal C(U)\). Consequently \(U\equiv V\)
  if and only if the two sentences have the same maximal normal form, if
  and only if they have the same minimal normal form.
\item If \(U\entails V\), then \(g_1(U)\ge g_1(V)\) and
  \(r_1(U)\le r_1(V)\); and \(0\le g_1(U)\le1\), with \(g_1(U)=0\)
  exactly when \(U\) is valid and \(g_1(U)=1\) exactly when \(U\) is
  contradictory.
\item For finite sets \(K,K^\ast\) of singular sentences whose
  conjunctions are equivalent, \(g_1(K)=g_1(K^\ast)\). If
  \(K\subseteq K^\ast\), then \(g_1(K)\le g_1(K^\ast)\).
\item If a structure \(\mathfrak M\) has ground assignment \(v\), then
  for every singular \(U\), \(\mathfrak M\models U\) if and only if
  \(v\in[U]\); in particular \(\mathfrak M\models D_w\) exactly when
  \(w=v\).
\end{enumerate}
For a general range measure \(r\), item~2 holds in the form
\(U\entails V\Rightarrow r(U)\le r(V)\) (\cref{thm:p15-power-range-laws}),
but its converse can fail when assignments have zero weight
(\cref{thm:p15-power-strength-boundary}).
\end{hbtheorem}

Item~4 is the exact sense of describing a state of the world as
completely as the language allows: a state description fixes the truth
value of every singular sentence. It does not fix the truth value of a
quantified sentence, because the universe of discourse is not restricted
to the named objects. One ground assignment can extend both to a model
of \(\forall x(P(x)\to Q(x))\) and to a model of its negation
(\cref{thm:p15-power-finite-procedure}). The reading of a maximal normal
form as a list of possible states of the world is therefore a reading
for singular sentences.

\section{The uniform measures and Carnap's objection}
\label{sec:p15-w02c-p15-029-uniform}

Hempel and Oppenheim introduce \(r_1(U)=l_U/m\), (9.4), in strict
analogy to \(g_1\), and note that \(r_1(U)+g_1(U)=1\) for every singular
\(U\), (9.5). They then report that Carnap has shown that \(r_1\) gives
the corresponding probability \(p_1\), defined from \(r_1\) by (9.3),
characteristics incompatible with the intended meaning of logical
probability; this is why Carnap, and Helmer together with Hempel and
Oppenheim, proposed other range measures
\cite[p.~170, n.~39]{HempelOppenheim1948Explanation}. The function
\(r_1\) gives every maximal sentence the same weight.

The place they cite is Carnap's 1945 discussion of the confirmation
function \(c_w\), based on the measure \(m_w\) that gives every state
description the same value \cite[pp.~80--81, (6)]{Carnap1945InductiveLogic}.
In a language with one primitive predicate \(P\) and \(101\) individuals,
let \(e\) be the conjunction of \(P(a_1),\ldots,P(a_{100})\) and let
\(h\) be \(P(a_{101})\). Then \(c_w(h,e)=1/2\), and the value stays
\(1/2\) if any or all of the conjuncts of \(e\) are replaced by their
negations. Carnap concludes that choosing \(c_w\) would amount to never
letting past experience influence expectations about the future
\cite[p.~81]{Carnap1945InductiveLogic}. On singular sentences \(m_w\) and
\(r_1\) are the same measure, so \(c_w=p_1=p_{r_1}\) there; in the 1950
presentation the same function is \(c^\dagger\) of
\cref{def:p15-carnap-star}. The objection extends to systematic power and
to every vocabulary, relational or not.

\begin{hbtheorem}{thm:p15-w02c-p15-029-uniform-independence}{proved=annex:p15-w02c-p15-029-uniform-independence}{Under the uniform measures, disjoint atoms are independent}
Let \(T\) and \(S\) be singular sentences such that no ground atom occurs
in both. Then
\[
 r_1(T\land S)=r_1(T)\,r_1(S),
\]
so \(p_{r_1}(T,S)=r_1(T)\) whenever \(r_1(S)>0\), and
\(s_{r_1}(T,S)=g_1(T)\) whenever \(g_1(S)>0\). In Carnap's example,
\(p_{r_1}(h,e)=s_{r_1}(h,e)=1/2\), and the same holds for every \(e'\)
obtained from \(e\) by negating conjuncts.

By contrast, let \(r_2=m^\ast\) be Carnap's measure of
\cref{def:p15-carnap-star} on the singular sentences of the language
with names \(a,b\) and one unary predicate \(P\). Its weights on
\(P(a)\land P(b)\), \(P(a)\land\neg P(b)\), \(\neg P(a)\land P(b)\),
\(\neg P(a)\land\neg P(b)\) are \(1/3,1/6,1/6,1/3\), and
\[
 p_{r_2}(P(b),P(a))=s_{r_2}(P(b),P(a))=\frac23,
 \qquad
 r_2(P(b))=g_{r_2}(P(b))=\frac12 .
\]
\end{hbtheorem}

Under \(r_1\), therefore, the systematic power of a theory about some
atoms, relative to data about other atoms, is just the theory's own
content \(g_1(T)\). By \cref{thm:p15-w02b-p15-025-increment} the data are
then irrelevant to the theory: systematic power above content is a
relevance quantity, and under \(r_1\) it vanishes for every pair with
disjoint atoms. Under \(r_2\) the excess is \(2/3-1/2=1/6>0\). A second
measure with positive excess, a latent mixture, is in
\cref{thm:p15-power-weight-choice}.

\section{The generalized measure and its evaluation}
\label{sec:p15-w02c-p15-029-evaluation}

\paragraph{Carnap's unequal weights.}
Hempel and Oppenheim describe Carnap's alternative as follows. A rule
assigns a specific weight, that is, a value of the range measure, to each
maximal sentence, and the weights are not the same for all maximal
sentences. The range measure of any other singular sentence is the sum of
the weights of the maximal sentences in its maximal normal form. Calling
the result \(r_2\), they write Carnap's probability as
\(p_2(T,S)=r_2(T\cdot S)/r_2(S)\) for singular \(T,S\), and note that
Carnap extends \(r_2\) and \(p_2\) by limiting processes to the cases in
which \(T\) and \(S\) are not both singular
\cite[p.~170]{HempelOppenheim1948Explanation}. They do not write out the
rule. In the 1945 paper the rule of this kind is \(m^\ast\): each
structure description receives the same value, and that value is divided
equally among the state descriptions isomorphic to one another
\cite[p.~80, (3)--(5)]{Carnap1945InductiveLogic}. Its 1950 form, with
its consequences, is \cref{def:p15-carnap-star} and
\cref{thm:p15-carnap-star-construction}. On singular sentences
\(r_2\) satisfies the conditions (9.7), as every positive weighting of
the assignments does (\cref{thm:p15-power-range-laws}).

\paragraph{The limit definition set aside.}
In note~38 Hempel and Oppenheim record that in Carnap's theory
\(p(T,S)\) is defined, in certain cases, as the limit that
\(r(T\cdot S)/r(S)\) approaches under specified conditions, and they
decline to consider this generalization of the schema (9.3)
\cite[p.~170, n.~38]{HempelOppenheim1948Explanation}. In the cited
place Carnap defines the measure and the confirmation function of the
infinite language as limits of their values in the finite languages
\(L_N\) as \(N\) grows. These definitions apply only where the limits
exist \cite[p.~75, (6)--(7)]{Carnap1945InductiveLogic}. That is the
material difference from (9.7), whose first condition requires \(r(U)\)
to be determined for every sentence \(U\). The handbook's treatment of
the infinite-system limits is in \cref{sec:p15-w02b-p15-029}.

\paragraph{The information-relative measure of note 40.}
The alternative of Helmer, Hempel and Oppenheim uses a range measure
\(r_I\) that depends in a specified way on the available empirical
information \(I\). The range measure of a sentence is therefore
determined only once a sentence \(I\) expressing that information is
given. From \(r_I\) a degree of confirmation \(dc\) is defined by a
formula similar to (9.3). Note~40 records three further points. First,
\(dc(T,S)\) is undefined in certain cases where \(S\) is generalized, as
McKinsey had pointed out. Second, \(dc\) does not satisfy all the
theorems of elementary probability theory, as discussed in Helmer and
Oppenheim's 1945 paper and in Hempel and Oppenheim's paper of the same
year. Hence the degree of confirmation of a theory on given evidence is
not a probability in the strict sense. Third, the definition of \(dc\)
has certain methodologically desirable features. It might therefore be of
interest to build a related concept of systematic power from \(r_I\), but
that question is not pursued in the paper
\cite[pp.~170--171, n.~40]{HempelOppenheim1948Explanation}. The
construction of \(r_I\) itself is given in those two 1945 papers, not in
the 1948 paper, and it is not presented in this handbook. What can be
stated exactly from the 1948 text alone is the following.

\begin{hblemma}{thm:p15-w02c-p15-029-strict-probability}{proved=annex:p15-w02c-p15-029-strict-probability}{A function that breaks an elementary law is not a range-measure probability}
Let \(c\) be a real-valued function defined on some pairs of sentences of
a fixed language. Suppose that, at some sentences \(T,T_1,T_2,S\), the
terms of one of the identities of \cref{thm:p15-power-eight-identities}
written with \(c\) in place of \(p_r\) are all defined, the side
condition of that identity holds, and the identity fails. Then there is
no range measure \(r\) on the language such that \(c(U,V)=p_r(U,V)\),
with \(r(V)>0\), at every pair \((U,V)\) occurring in those terms.
Likewise there is no range measure \(r\) such that
\(c(U,V)=s_r(\neg U,\neg V)\), with \(g(\neg V)>0\), at every such pair:
\(c\) is not the systematic-power function of any range measure read at
the negated arguments.
\end{hblemma}

So the remark that \(dc\) is not a probability in the strict sense
implies that no single range measure satisfying (9.7) yields \(dc\)
through (9.10) at the pairs where it fails a law. The results of
\cref{sec:p15-systematic-power} and of
\cref{sec:p15-w02b-p15-025-comparison} therefore cannot be applied to
\(dc\), or to a systematic-power function built from \(r_I\), without
first checking that \(r_I\), for the fixed \(I\), satisfies (9.7).
Point~2 of \cref{def:p15-w02b-p15-025-limits} records the same
restriction from the side of the comparison with relevance measures.

\paragraph{The extended domain.}
The function \(s(T,S)=g(T\vee S)/g(S)\) of (9.9) is defined for every
sentence \(T\) and every sentence \(S\) with \(g(S)\neq0\). The
definition of section~8, by contrast, was restricted to singular data
that are not formally true \cite[p.~172]{HempelOppenheim1948Explanation}.
The next theorem states exactly how the two definitions are related when
a range measure on the whole language agrees with \(r_1\) on singular
sentences.

\begin{hbtheorem}{thm:p15-w02c-p15-029-extension}{proved=annex:p15-w02c-p15-029-extension}{Where the generalized measure agrees with the original one, and where it does not}
Let \(r\) be a range measure on all sentences of the language with
\(r(U)=r_1(U)\) for every singular \(U\), and let \(g=1-r\). Let \(S\)
be singular and not valid, and let \(T^\flat\) be the singular
consequence projection of \cref{thm:p15-power-uniform-formula}.
\begin{enumerate}
\item For every sentence \(T\),
  \[
   s_r(T,S)=s_0(T,\{S\})+\frac{r(\neg T\land T^\flat\land\neg S)}{g_1(S)}
   \ \ge\ s_0(T,\{S\}).
  \]
\item If \(T\) is singular, then \(s_r(T,S)=s_0(T,\{S\})\).
\item In \(L_1\), let \(T=\forall x(P(x)\to Q(x))\) and
  \(S=P(a)\land Q(a)\). There are range measures \(r'\) and \(r''\) on
  all sentences of \(L_1\), each agreeing with \(r_1\) on singular
  sentences, with
  \[
   s_{r'}(T,S)=\tfrac13=s_0(T,\{P(a),Q(a)\}),\qquad
   s_{r''}(T,S)=\tfrac23 .
  \]
\item Under \(r'\), the non-singular sentence \(\forall x\,P(x)\) has
  \(g_{r'}(\forall x\,P(x))=1/2\), so \(s_{r'}(U,\forall x\,P(x))\) is
  defined for every \(U\), although the definition of section~8 does
  not apply to these data.
\end{enumerate}
\end{hbtheorem}

\paragraph{The source's evaluation.}
Hempel and Oppenheim close Part~IV by comparing the two stages. They
began with the optimum ratio of the data derivable from the remainder by
means of the theory. Working that idea out produced a broader notion of
systematic power, and this broader notion turned out to be dual to
logical probability. In their judgement this extension gives a simpler and more
comprehensive theory than the initial definition could have supported
\cite[p.~173]{HempelOppenheim1948Explanation}. They add that the theory
of systematic power, in both versions, is purely formal, as is the theory
of logical probability. A significant application of either in
epistemology or the methodology of science requires first settling basic
questions about how the language of science is logically built and how
its concepts are to be interpreted. Two such problems are the requirement of
purely qualitative primitives and the choice of an adequate range
measure among infinitely many formal possibilities
\cite[p.~173]{HempelOppenheim1948Explanation}. For the complexity of
these issues they cite, in note~41, Goodman's ``A query on
confirmation'', ``The problem of counterfactual conditionals'' and ``On
infirmities of confirmation theory'', and Carnap's ``On the application
of inductive logic'', and point back to their own notes~21 and~25
\cite[p.~173, n.~41]{HempelOppenheim1948Explanation}.

Both comparative claims have exact counterparts. The theory is
comprehensive in the sense of \cref{thm:p15-w02c-p15-029-extension}: it
reproduces the section~8 values for singular theories and extends the
measure to non-singular data. It is simpler in that the relation to
probability is the one-line identity (9.11),
\(s(T,S)=p(\neg T,\neg S)\) (\cref{thm:p15-power-probability-dual}), and
every law about \(p\) transfers to \(s\) and back by the Dualism Theorem
(\cref{thm:p15-power-dualism} and \cref{thm:p15-w02c-p15-029-dualism-converse}).
The section~8 definition instead needs the combinatorial optimal-split
argument of \cref{thm:p15-power-optimal-split}. The comprehensiveness
has a price, which item~3 of \cref{thm:p15-w02c-p15-029-extension} makes
visible. For a quantified theory, the generalized value depends on how
the range measure weighs non-singular sentences, which (9.7) leaves
open. Fixing \(r_1\) on singular sentences does not fix it. This is the
second open problem of the closing paragraph in a concrete case.

\section{The converse half of the Dualism Theorem}
\label{sec:p15-w02c-p15-029-dualism}

The Dualism Theorem (9.13) states two transformations. Replacing
\(p\) by \(s\) and exchanging conjunction and disjunction turns any
demonstrable general equality or inequality about \(p\) into one about
\(s\). The same exchange, with \(s\) replaced by \(p\), turns any such
theorem about \(s\) into one about \(p\). Hempel and Oppenheim state the
theorem informally and say that it can be proved from (9.11)
\cite[pp.~172--173]{HempelOppenheim1948Explanation}. The forward half,
in the precise form of \cref{thm:p15-power-dualism}, is proved in
\cref{annex:p15.power-dualism}. The converse half requires its own
identity between terms, its own transfer of definedness conditions, and
a substitution of negated sentences into a valid statement about \(s\).

\begin{hbtheorem}{thm:p15-w02c-p15-029-dualism-converse}{proved=annex:p15-w02c-p15-029-converse}{The converse transformation of source Theorem 9.13}
Let \(A_1,\ldots,A_m\) be sentence metavariables and \(D\) the operation
of \cref{def:p15-power-dual-syntax}. For Boolean expressions \(B\) and
sentences \(U_1,\ldots,U_m\), write \(B(U)\) for
\(B(U_1,\ldots,U_m)\) and \(B(\neg U)\) for
\(B(\neg U_1,\ldots,\neg U_m)\).
\begin{enumerate}
\item \(D(D(B))=B\) for every Boolean expression \(B\).
\item For every range measure \(r\), all Boolean expressions \(B,C\) and
  all sentences \(U_1,\ldots,U_m\),
  \[
   r\bigl(D(C)(U)\bigr)=g\bigl(C(\neg U)\bigr),
  \]
  so \(p_r(D(B)(U),D(C)(U))\) is defined exactly when
  \(s_r(B(\neg U),C(\neg U))\) is, and then
  \[
   p_r\bigl(D(B)(U),D(C)(U)\bigr)=s_r\bigl(B(\neg U),C(\neg U)\bigr).
  \]
\item Call a statement \(\Psi\) built from terms \(s_r(B_j,C_j)\),
  numbers, numerical operations, equalities and inequalities, and
  possibly hypotheses of the form \(B\entails C\), \emph{universally
  valid} if it is true for every language of
  \cref{def:p15-power-language}, every range measure on it and every
  choice of sentences for the metavariables at which all its terms are
  defined and all its hypotheses hold. Let \(\Psi^D\) be obtained by
  replacing each \(s_r(B_j,C_j)\) by \(p_r(D(B_j),D(C_j))\) and each
  hypothesis \(B\entails C\) by \(D(C)\entails D(B)\). If \(\Psi\) is
  universally valid, so is \(\Psi^D\), with \(p\)-statements understood
  in the same way.
\item Applying the forward transformation of \cref{thm:p15-power-dualism}
  to \(\Psi^D\) returns \(\Psi\). The two transformations are inverse to
  each other.
\end{enumerate}
For example, the converse transformation sends the power product law of
(9.12), defined when \(g(S)>0\) and \(g(T_1\lor S)>0\), to the
probability product law, defined when \(r(S)>0\) and
\(r(T_1\land S)>0\):
\[
 s_r(T_1\lor T_2,S)=s_r(T_1,S)\,s_r(T_2,T_1\lor S)
\]
becomes
\[
 p_r(T_1\land T_2,S)=p_r(T_1,S)\,p_r(T_2,T_1\land S).
\]
It sends ``if \(T\entails U\) then
\(s_r(T,S)\ge s_r(U,S)\)'' (\cref{thm:p15-power-strength-boundary}) to
``if \(U\entails T\) then \(p_r(T,S)\ge p_r(U,S)\)''.
\end{hbtheorem}

\begin{hbwarning}{def:p15-w02c-p15-029-naive-converse}{Replacing letters is not the converse transformation}
The truth constants must be exchanged together with the connectives.
The statement \(s_r(\bot,S)=1\) holds whenever \(g(S)>0\)
(\cref{thm:p15-power-strength-boundary}). Replacing \(s\) by \(p\) while
keeping \(\bot\) gives \(p_r(\bot,S)=1\), which is false whenever
\(r(S)>0\), because the left side is \(0\). The converse transformation
gives \(p_r(\top,S)=1\), which is true. Likewise a hypothesis must be
reversed, not copied: copying \(T\entails U\) into the probability
statement would give ``if \(T\entails U\) then
\(p_r(T,S)\ge p_r(U,S)\)'', which fails for \(T=\bot\), \(U=\top\) and
any \(S\) with \(r(S)>0\).
\end{hbwarning}

\section{Limitations, relationships and sources}
\label{sec:p15-w02c-p15-029-limits}

\paragraph{Limitations.}
The source reports in this chapter state what Hempel and Oppenheim claim;
their judgements of simplicity and comprehensiveness, and Carnap's
judgement on \(c_w\), are evaluations, not theorems. The exact
counterparts given here are proved for the stated measures and languages
only. \Cref{thm:p15-w02c-p15-029-uniform-independence} concerns \(r_1\)
and one instance of \(r_2\); it does not describe Carnap's \(m^\ast\) in
general, which is treated in \cref{sec:p15-carnap-star}.
\Cref{thm:p15-w02c-p15-029-extension} assumes a range measure on the
whole language that agrees with \(r_1\) on singular sentences. Carnap's
limit definitions need not give such a measure, because their limits
need not exist for every sentence. The measure \(r_I\) and the function
\(dc\) of note~40 are reported, not constructed.

\paragraph{Relationships.}
The general constructions are in \cref{sec:p15-systematic-power}, and
the model-language computations are in
\cref{sec:p15-w02c-p15-028-model-languages}. The relevance reading of
systematic power above content, used after
\cref{thm:p15-w02c-p15-029-uniform-independence}, is
\cref{thm:p15-w02b-p15-025-increment}. Carnap's \(m^\ast\), \(c^\ast\)
and the equal-state function \(c^\dagger\) are in
\cref{sec:p15-carnap-star}. The prediction criterion, of which
systematic power is a quantitative counterpart, is
\cref{def:p15-hempel-prediction}.

\paragraph{Sources.}
Hempel and Oppenheim 1948, Part~IV: the intuitive characterization and
the ratio, p.~164; potential data, the theory argument, note~33 and the
percentage normalization, p.~165; maximal sentences, duality and the
content indicator, p.~168 with note~36; the range indicator, possible
states of the world, exclusive ranges, note~37, inverse variation, (9.1)
and \(g_1\) of a class, pp.~168--169; (9.4), (9.5), Carnap's objection
with note~39, \(r_2\), \(p_2\) and note~38, p.~170; note~40,
pp.~170--171; the domain of (9.9), p.~172; (9.13) and the closing
evaluation with note~41, p.~173
\cite[pp.~164--173]{HempelOppenheim1948Explanation}. Carnap 1945: the
limit definitions, p.~75, (6)--(7); \(m^\ast\), p.~80, (3)--(5); \(m_w\),
\(c_w\) and the objection to them, pp.~80--81
\cite{Carnap1945InductiveLogic}.

%% toc: P15 164 164.12
\chapter{Romeijn on Inductive Logic and Statistics}
\label{sec:p15-w02c-p15-029-romeijn}

Jan-Willem Romeijn's chapter ``Inductive Logic and Statistics'' in the
tenth volume of the \emph{Handbook of the History of Logic} asks whether
the inferences of a probabilistic inductive logic can represent the
procedures of statistics \cite[pp.~625--650]{Romeijn2011InductiveLogicStatistics}.

\paragraph{The two disciplines and how they relate.}
Romeijn opens by defining both disciplines
\cite[p.~625]{Romeijn2011InductiveLogicStatistics}. An inductive logic
is a system of inference relating propositions about data to
propositions that reach past the data; his examples of the latter are
predictions about future data and general conclusions drawn from the
data. Statistics is a mathematical discipline whose procedures take
sample data to results about a population. The results he lists are
decisions to reject or to accept a hypothesis about the population,
the fixing of probability assignments over such hypotheses, and
predictions about future samples, and he leaves the list open. What
relates the two, on his account, is that each is a calculus leading
from the data in hand to propositions or results that go beyond those
data \cite[p.~625]{Romeijn2011InductiveLogicStatistics}. These are
informal characterizations; their formal counterparts in the chapter
are the observation algebra of \cref{sec:p15-w02c-p15-029-rj-data} and
the two templates of \cref{def:p15-w02c-p15-029-rj-templates}.

\paragraph{Why the two developed apart.}
In spite of this common character, Romeijn says, the two disciplines
developed more or less separately, and he gives two reasons
\cite[p.~625]{Romeijn2011InductiveLogicStatistics}. The first, which he
says explains the separation only in part, is that there are objections
to regarding statistics as inferential at all, and above all to
regarding classical statistical procedures so. The second, which he
regards as the more important, is that the Carnapian programme has
dominated inductive logic, and that statisticians may not have seen in
Carnapian inductive logic a discipline much like their own. He traces
this second reason, and much of the mutual misunderstanding, to the
different roles the two programmes give to hypotheses: statistical
hypotheses and models are absent from Carnap's logic, while most
statistical procedures begin and end with them
\cite[p.~625]{Romeijn2011InductiveLogicStatistics}. Both reasons are
historical claims, and he marks the weighting between them and the
diagnosis of the misunderstanding as his own opinion rather than as a
documented finding. On the next page he carries the diagnosis one step
further back, to the logical-empiricist roots of inductive logic (see
the scope paragraph below).

\paragraph{Thesis and aims.}
The thesis is that Carnapian inductive logic can be extended to
inference over statistical hypotheses, and that an inductive logic
whose language contains such hypotheses can capture Bayesian
statistics exactly and Neyman--Pearson testing and Fisher's estimation
in part \cite[pp.~625--626]{Romeijn2011InductiveLogicStatistics}.
Romeijn expects three gains from bringing the disciplines together:
reading statistical procedures as inferences may clarify their
presuppositions and foundations; insights from inductive logic may be
used to improve statistics; and inductive logicians may come to see the
bearing of their discipline on the sciences
\cite[p.~625]{Romeijn2011InductiveLogicStatistics}.
This chapter presents the whole of Romeijn's chapter in its own
notation: the observation language and sample space, the two templates
for an inductive step, Carnap's symmetry principles and
\(\lambda\)-rules, Bayesian inference over multinomial hypotheses and
its link to Carnap through de Finetti's theorem, Neyman--Pearson
testing and its Bayesian imitation, maximum likelihood estimation with
confidence and credal intervals, and the fiducial argument with its
functional-model presuppositions.

\paragraph{Purpose, inputs and conclusions.}
The input of every construction is a sample space of observation
sequences with a family of probability functions on it, together with
whatever premises restrict that family. The conclusions are of two
kinds and are kept apart below. Romeijn's conclusions are a
classification of inferential steps (ampliative or not), the claim that
Carnapian logic and Bayesian statistics are non-ampliative and
correspond through de Finetti's representation, and the claim that
classical procedures are best read as ampliative but admit partial
non-ampliative imitations. The editorial conclusions are the theorems
of this chapter, each proved in the annex, which verify the chapter's
calculations, give its general claims exact hypotheses, and mark where
a claim fails without them.

\paragraph{Claim types.}
Three kinds of claim are distinguished throughout. A
\emph{mathematical} claim is a statement about probability functions
that is true or false under stated hypotheses; the source's numbered
equations are of this kind. A \emph{philosophical} claim concerns the
nature, interpretation or justification of inference, such as the
claim that statistical hypotheses are ``elusive'' for a logical
empiricist. A \emph{historical} claim reports who held or proved what,
and when. The chapter reports no empirical study. Its two descriptive
remarks about practice, that users often read a significance level as
the probability that the null hypothesis is true and that the sciences
persist in using statistical hypotheses, are given without a task,
population, design or measure \cite[pp.~636, 640]{Romeijn2011InductiveLogicStatistics};
they are reported here as the author's observations and no conclusion
is drawn from them.

\paragraph{Scope the chapter sets for itself.}
Romeijn treats positions from both disciplines systematically rather
than as historical entities. Its connection to history is indirect: he
traces the small common past of the two fields to the absence of
statistical hypotheses from inductive logic, and that absence to the
logical-empiricist roots of inductive logic \cite[p.~626]{Romeijn2011InductiveLogicStatistics}.
He says the discussion of statistics is short, that many statistical
procedures are not treated, and that the many Carnapian systems are not
discussed in detail, referring for these to chapters 9 and 10 of the
volume and to the philosophy-of-statistics volume edited by
Bandyopadhyay and Forster \cite[p.~626]{Romeijn2011InductiveLogicStatistics}.

\section{Observations, the observation language and the sample space}
\label{sec:p15-w02c-p15-029-rj-data}

An inductive inference starts from propositions about data and ends in
propositions that say more than the data classically entail: from the
fact that every pear seen so far was green, that the next few pears will
be green, or that all pears are green \cite[p.~626]{Romeijn2011InductiveLogicStatistics}.
Romeijn restricts the data in three ways. Only propositions about
empirical facts count, so that ``pears are healthy'' is excluded. Only
results of observing a designated kind of empirical fact count, such as
the colour of individual pears. And each result belongs to a fixed
partition of mutually exclusive properties, typically finite and always
countable; data are always a finite sequence of such results, although
there may be infinitely many possible results \cite[pp.~626--627]{Romeijn2011InductiveLogicStatistics}.

\begin{hbdefinition}{def:p15-w02c-p15-029-rj-observations}{Observation language and sample space}
Properties are numbered \(k\in K=\{0,1,\ldots,n-1\}\). The observation
of individual \(i\) is the propositional variable \(Q_i\); its valuation
\(q^k_i\) is the sentence that individual \(i\) has property \(k\). The
first \(t\) results form the variable \(S_t\), with valuation
\(s^{k_1\ldots k_t}\), abbreviated \(s_t\). The expressions \(s_t\) and
\(q^k_i\) make up the \emph{observation language}.

Let \(K^{\Omega}\) be the set of all infinite sequences of properties,
with elements \(\omega\) and \(i\)-th entry \(\omega(i)\). Each sentence
is identified with a set of sequences:
\[
 q^k_i=\{\omega\in K^{\Omega}:\ \omega(i)=k\},\qquad
 s^{k_1\ldots k_t}=\bigcap_{i=1}^{t} q^{k_i}_i .
\]
The \emph{sample space} or \emph{observation algebra} \(\mathcal Q\) is
the algebra built from the sets \(q^k_i\) by the Boolean operations.
Entailment between sentences becomes inclusion between sets. If
\(\mathcal Q\) is the \(\sigma\)-algebra generated by the observations,
it also contains the set \(\{\forall i: q^1_i\}\) of sequences in which
every individual has property \(1\)
\cite[pp.~627--628 and n.~1]{Romeijn2011InductiveLogicStatistics}.
\end{hbdefinition}

In the running example the properties are the colours red, green and
yellow, written \(0\), \(1\) and \(2\), so \(n=3\), and \(s^{012}\) says
that the first three pears were red, green and yellow; then
\(s^{012}\subset q^1_2\), the set form of \(s^{012}\to q^1_2\)
\cite[p.~627]{Romeijn2011InductiveLogicStatistics}. Romeijn stresses
the generality of the construction: it excludes only a continuum of
individuals and a continuum of properties, so subjects with test scores
or companies with stock prices fit as well as pears with colours, and
it has the basic structure of most statistical applications and almost
all applications in inductive logic \cite[pp.~627--628]{Romeijn2011InductiveLogicStatistics}.
He notes that probability is strictly a measure on sets, but that its
arguments may equally be read as sentences
\cite[p.~627]{Romeijn2011InductiveLogicStatistics}. The symbol
\(\omega\) has a second, unrelated use in the fiducial section, where it
names an error term (\cref{def:p15-w02c-p15-029-rj-functional-model}).

\section{Inductive inference: two templates}
\label{sec:p15-w02c-p15-029-rj-inference}

From the sample space alone nothing inductive follows. After three red
pears \(s^{000}\), each of the sets \(s^{000k}=s^{000}\cap q^k_4\) for
\(k=0,1,2\) is a nonempty member of the space, so no \(q^k_4\) is
derivable from \(s^{000}\) \cite[p.~628]{Romeijn2011InductiveLogicStatistics}.
Probabilistic inductive logic adds a probability function
\(P:\mathcal Q\to[0,1]\). Every possible next pear \(q^k_{t+1}\) and
every sequence \(s^{k_1\ldots k_t}\) receives a probability; predictions
are conditional probabilities \(P(q^k_{t+1}\mid s^{k_1\ldots k_t})\), and
the probability that all pears are green is \(P(\{\forall i:q^1_i\})\).
All such inferences are fixed by the full function \(P\), so the central
question of any inductive procedure is how \(P\) is to be determined
relative to the data in hand; statistics and inductive logic answer it
differently \cite[p.~628]{Romeijn2011InductiveLogicStatistics}.

Romeijn reads probability functions as models. A model of classical
two-valued logic is a complete truth valuation, which, through the
correspondence of language and algebra, is a \(\{0,1\}\)-valued function
on the algebra. A probability function on \(\mathcal Q\) is a model in
the same sense, with values in \([0,1]\) and subject to the probability
axioms. Because statisticians use ``model'' for sets of statistical
hypotheses, he calls a complete probability function on the algebra a
\emph{statistical hypothesis} \(h\), with probability function \(P_h\),
also called a distribution over a population
\cite[pp.~628--629]{Romeijn2011InductiveLogicStatistics}.

\begin{hbdefinition}{def:p15-w02c-p15-029-rj-templates}{Ampliative and non-ampliative steps}
Romeijn distinguishes two ways of making the inductive step precise
\cite[p.~629]{Romeijn2011InductiveLogicStatistics}.
\begin{enumerate}
\item \emph{Ampliative.} The step maps a sample \(s_t\), together with a
  set of statistical hypotheses each with its function \(P_h\), to a
  smaller set of hypotheses or to a single hypothesis \(h^\star\) with
  \(P_{h^\star}\). It is ampliative because the restriction effected by
  the data (the conclusion) is often stronger than what the data and the
  initial set of functions (the premises) deductively entail.
\item \emph{Non-ampliative.} Premises are restrictions on probability
  assignments over the sample space; each determines the set of
  probability functions satisfying it, and a conclusion likewise
  determines a set. The inference is \emph{valid} when the set of
  functions satisfying the premises is contained in the set satisfying
  the conclusion, exactly as classical validity is inclusion of the
  premises' models in the conclusion's models. A valid step does not
  amplify its premises.
\end{enumerate}
The second template is the one Romeijn associates with what has
traditionally been called inductive logic, with the probabilistic logic
of Hailperin (1996) and of Haenni and others (2009, especially
section~2), and with the Bayesian logic advocated by Howson (2003)
\cite[p.~629]{Romeijn2011InductiveLogicStatistics}.
\end{hbdefinition}

The chapter gives two examples of a valid non-ampliative step: fixing
\(P(q^0_1)=\tfrac12\) and \(P(q^1_1)=\tfrac13\) entails
\(P(q^2_1)=\tfrac16\); and the posterior \(P(h\mid s_t)\) is obtained by
Bayes's theorem from \(P(h)\), \(P(s_t)\) and \(P(s_t\mid h)\)
\cite[p.~629]{Romeijn2011InductiveLogicStatistics}. The first example
presupposes the three-colour partition. The next theorem states both with
the side conditions under which the inclusion of model sets holds.

\begin{hbtheorem}{thm:p15-w02c-p15-029-rj-entailment}{proved=annex:p15-w02c-p15-029-rj-entailment}{The two examples of non-ampliative entailment}
\begin{enumerate}
\item Let \(n=3\). Every finitely additive probability function \(P\) on
  \(\mathcal Q\) with \(P(q^0_1)=\tfrac12\) and \(P(q^1_1)=\tfrac13\)
  satisfies \(P(q^2_1)=\tfrac16\), and such functions exist. So the set
  of models of the two premises is a nonempty subset of the set of
  models of the conclusion.
\item Let \(P\) be a probability function on an algebra containing a
  hypothesis \(h\) and a sample \(s_t\). If \(P(h)=a>0\),
  \(P(s_t\mid h)=b\) and \(P(s_t)=c>0\), then \(P(h\mid s_t)=ab/c\).
  Hence every model of these three premises is a model of
  \(P(h\mid s_t)=ab/c\).
\end{enumerate}
\end{hbtheorem}

The chapter concentrates on the non-ampliative template because Carnapian
logic fits it most easily, so that reading statistical procedures this
way makes them accessible to inductive-logical analysis. He disclaims
two stronger aims: he does not claim to expose the real nature of the
statistical procedures, nor to supply indubitable norms for statistical
inference; the aim is clarification, enrichment, and research on
inductive logic directed at statistical problems
\cite[p.~630]{Romeijn2011InductiveLogicStatistics}.

\section{Carnapian logics}
\label{sec:p15-w02c-p15-029-rj-carnap}

Romeijn credits Carnap's logic of induction (1950, 1952) with the best
claim to the title of inductive logic proper and calls it the first
systematic study of probabilistic prediction from data
\cite[p.~630]{Romeijn2011InductiveLogicStatistics}. Its central notion
is logical probability: a probability assignment derived from
symmetries of the observation language. With three mutually exclusive
colours and no further knowledge, nothing singles out one colour, which
suggests \(P(q^0_i)=P(q^1_i)=P(q^2_i)=\tfrac13\). The aim is a unique
function, or a strongly restricted set, fixed by such symmetries
\cite[p.~630]{Romeijn2011InductiveLogicStatistics}.

Predictive properties restrict the set further. If yellow pears are
taken to be more like green pears than red ones, finding a yellow pear
should lower the probability of red markedly and that of green much
less; in ratio form,
\[
 \frac{P(q^1_{t+1}\mid s_{t-1}\cap q^2_t)}{P(q^0_{t+1}\mid s_{t-1}\cap q^2_t)}
 >
 \frac{P(q^1_{t+1}\mid s_{t-1})}{P(q^0_{t+1}\mid s_{t-1})}.
\]
Romeijn refers for such analogy properties to Festa (1996), Maher
(2000), his own paper of 2006, and Paris and Waterhouse (2008)
\cite[p.~630]{Romeijn2011InductiveLogicStatistics}.

\begin{hbdefinition}{def:p15-w02c-p15-029-rj-lambda}{Carnap's symmetries, restricted relevance and \texorpdfstring{\(\lambda\)}{lambda}-rules}
As Romeijn presents it, the inductive logic discussed at length in
Carnap (1952) uses a probability function satisfying
\begin{align*}
 &\text{(1)}\quad P(q^k_i)=P(q^{k'}_i),\\
 &\text{(2)}\quad P(s^{k_1\ldots k_i\ldots k_t})=P(s^{k_i\ldots k_1\ldots k_t}),
\end{align*}
for all \(i,t,k,k'\) and all \(k_i\) with \(1\le i\le t\). Condition (2)
is \emph{exchangeability}: the order of observations does not affect
their probability. The logic uses the particular version called
\emph{restricted relevance},
\[
 \text{(3)}\quad P(q^k_{t+1}\mid s_t)=f(t_k,t),
\]
where \(t_k\) is the number of earlier instances \(q^k_i\) in \(s_t\) and
\(t\) the number of observations. From these Romeijn states that one
can derive
\[
 \text{(4)}\quad P(q^k_{t+1}\mid s_t)=\frac{t_k+\frac{\lambda}{n}}{t+\lambda},
\]
with \(n\) the number of values of \(k\) and a freely chosen parameter
\(0\le\lambda<\infty\). Predictive assignments of form (4) are
\emph{Carnapian \(\lambda\)-rules}. Their most important feature is
\emph{instantial relevance},
\[
 \text{(5)}\quad P(q^k_{t+1}\mid s_{t-1}\cap q^k_t)>P(q^k_{t+1}\mid s_{t-1}):
\]
an occurrence of \(q^k_t\) raises the probability of \(q^k_{t+1}\)
\cite[p.~631]{Romeijn2011InductiveLogicStatistics}.
\end{hbdefinition}

Romeijn reports that deriving this typically inductive effect from
symmetries alone was counted a success for Carnap: by justifying the
symmetries independently, Carnap in effect offered a justification of
induction and an answer to Hume. Footnote~2 adds the history: as
Zabell (1982) recounts, Johnson (1932) and de Finetti (1937) had
already connected exchangeability with predictive properties, but the
link to Hume's problem is Carnap's, who motivated properties such as
(4) through logical probability, whereas for the subjectivist de Finetti
they had no objective ground \cite[p.~631 and n.~2]{Romeijn2011InductiveLogicStatistics}.
That justification claim is philosophical and is reported, not
endorsed. The mathematical part, (5), holds for the \(\lambda\)-rules
under exact hypotheses that exclude the endpoint \(\lambda=0\).

\begin{hbtheorem}{thm:p15-w02c-p15-029-rj-instantial}{proved=annex:p15-w02c-p15-029-rj-instantial}{Instantial relevance for \texorpdfstring{\(\lambda\)}{lambda}-rules, and its failure at \texorpdfstring{\(\lambda=0\)}{lambda = 0}}
Let \(n\ge2\) and let \(P\) be a probability function on \(\mathcal Q\)
whose predictive probabilities satisfy (4) for every \(t\ge0\) and every
\(s_t\) with \(P(s_t)>0\).
\begin{enumerate}
\item If \(0<\lambda<\infty\), then \(P(s_t)>0\) for every finite
  sequence \(s_t\), and for every \(t\ge1\), every \(s_{t-1}\) with
  \(t'_k\) instances of \(k\),
  \[
   P(q^k_{t+1}\mid s_{t-1})=\frac{t'_k+\frac{\lambda}{n}}{t-1+\lambda}
   <\frac{t'_k+1+\frac{\lambda}{n}}{t+\lambda}
   =P(q^k_{t+1}\mid s_{t-1}\cap q^k_t).
  \]
  So instantial relevance (5) holds for every \(k\), \(t\) and
  \(s_{t-1}\).
\item If \(\lambda=0\), (4) does not define the first prediction
  (\(t=0\) gives \(0/0\)). Supplement it by any first-step values
  \(P(q^k_1)\). For \(t\ge2\) and a sequence \(s_{t-1}\) all of whose
  entries are \(k\), with \(P(s_{t-1})>0\), both sides of (5) equal
  \(1\), so the strict inequality fails; and if \(s_{t-1}\) has
  positive probability and no entry \(k\), then
  \(P(s_{t-1}\cap q^k_t)=0\) and the left side of (5) is undefined.
\end{enumerate}
\end{hbtheorem}

\begin{proofboundary}{def:p15-w02c-p15-029-rj-lambda-boundary}{What is and is not derived about (4)}
The chapter states that (1)--(3) yield (4) and does not prove it; he
refers to Carnap (1952) \cite[p.~631]{Romeijn2011InductiveLogicStatistics}.
That derivation is not reproduced here. The handbook's account of
Johnson's sufficientness postulate in \cref{sec:p15-carnap-history}
records that the characterization needs at least three categories and
that with two categories the postulate alone is vacuous; Romeijn's
statement carries no such qualification. Conversely, for
\(0<\lambda<\infty\) an exchangeable probability function on the
algebra generated by the sets \(q^k_i\) with predictive probabilities
(4) exists, so the hypotheses of
\cref{thm:p15-w02c-p15-029-rj-instantial} are satisfiable: take the
masses of \cref{def:p15-carnap-history-family} with \(k=n\) categories
and every \(\alpha_i=\lambda/n\); by
\cref{thm:p15-carnap-history-family} they are positive, normalized,
exchangeable and consistent under deleting the last observation, and
their predictive probability is \((t_k+\lambda/n)/(t+\lambda)\).
\end{proofboundary}

\paragraph{Carnapian logic as non-ampliative inference.}
Romeijn contrasts the outlook of Carnap's logic with the templates of
\cref{sec:p15-w02c-p15-029-rj-inference}. Those start from a set of
probability functions, or hypotheses, and restrict it or draw
consequences on the basis of data. Carnap starts from the sample space
and lets symmetry principles, jointly with predictive properties,
single out a set of probability functions, as truth tables restrict truth valuations,
though not to a single function, since \(\lambda\) stays free. Seen
from statistics, Carnap thereby motivates from logical principles the
choice of a particular set of hypotheses
\cite[pp.~631--632]{Romeijn2011InductiveLogicStatistics}. Setting
logical probability aside, the inferential step fits the non-ampliative
template: symmetries such as (1) and (2) act as premises restricting the
probability functions, and conclusions, chiefly probabilities
conditional on samples \(s_t\), are derived by the probability axioms
alone \cite[p.~632]{Romeijn2011InductiveLogicStatistics}.

\paragraph{Why hypotheses are absent.}
Both programmes work with sets of probability functions over the sample
space, but statistics frames them as statistical hypotheses and
Carnapian logic does not; it characterizes them by symmetries and
predictive properties. Romeijn locates the reason in logical
empiricism: symmetries concern the empirical predicates of the
language, and predictive properties concern what the probability
functions show on finite data, whereas statistical hypotheses concern
chances, cannot be stated as finite combinations of empirical
predicates, and show up, if at all, only as limiting relative
frequencies as the data grow without bound
\cite[p.~632]{Romeijn2011InductiveLogicStatistics}. He names Hintikka's
systems as the notable exception to the inductive logicians' avoidance
of hypotheses and refers to chapters 9 and 10 of the volume
\cite[p.~632]{Romeijn2011InductiveLogicStatistics}. These are
philosophical and historical claims.

\section{Bayesian statistics}
\label{sec:p15-w02c-p15-029-rj-bayes}

Bayesian statistics is the first procedure Romeijn examines. Its mark is
that probability assignments range over statistical hypotheses as well
as data \cite[pp.~632--633]{Romeijn2011InductiveLogicStatistics}.

\begin{hbdefinition}{def:p15-w02c-p15-029-rj-bayes-def}{Bayesian statistical inference (Romeijn's Definition 1)}
Let \(\mathcal H\) be the space of statistical hypotheses \(h_\theta\),
\(\theta\in\Theta\), the space of parameter values, and \(\mathcal Q\)
the sample space; the functions \(P\) are probability assignments over
\(\mathcal H\times\mathcal Q\), so \(P(s_t\mid h_\theta)\) ranges over
all of \(\mathcal Q\). Assume the prior \(P(h_\theta)\) and the
likelihoods \(P(s_t\mid h_\theta)\). Bayes's theorem gives
\[
 \text{(6)}\quad P(h_\theta\mid s_t)=\frac{P(h_\theta)P(s_t\mid h_\theta)}{P(s_t)},
\]
and Bayesian statistics outputs this posterior. The probability of the
data is not presupposed; it follows from prior and likelihoods by the
law of total probability,
\(P(s_t)=\int_\Theta P(h_\theta)P(s_t\mid h_\theta)\,d\theta\)
\cite[p.~633]{Romeijn2011InductiveLogicStatistics}.
\end{hbdefinition}

Estimates and measures of their accuracy are derived from the
posterior; Romeijn refers to Barnett (1999) and Press (2003). When only
two hypotheses are compared, the posterior ratio is
\[
 \frac{P(h_\theta\mid s_t)}{P(h_{\theta'}\mid s_t)}
 =\frac{P(h_\theta)P(s_t\mid h_\theta)}{P(h_{\theta'})P(s_t\mid h_{\theta'})},
\]
and with equal priors the ratio of likelihoods, which he calls the
Bayes factor, compares them \cite[p.~633]{Romeijn2011InductiveLogicStatistics}.

\paragraph{Emma's farm.}
Emma's pears are red, \(q^0_i\), or green, \(q^1_i\), in an unknown
ratio. The multinomial hypotheses \(h_\theta\), \(\theta\in[0,1]\), are
\[
 \text{(7)}\quad P_{h_\theta}(q^1_t\mid s_{t-1})=\theta,\qquad
 P_{h_\theta}(q^0_t\mid s_{t-1})=1-\theta:
\]
whatever was seen before, a randomly drawn pear is green with
probability \(\theta\); the induced distribution is called Bernoulli or
multinomial. Instead of choosing among hypotheses directly, the
Bayesian puts a distribution over them that expresses epistemic
uncertainty, for instance the Beta distribution
\[
 \text{(8)}\quad P(h_\theta)=\mathrm{Norm}\times\theta^{\lambda/2-1}(1-\theta)^{\lambda/2-1},
 \qquad\theta\in\Theta=[0,1],
\]
uniform for \(\lambda=2\). With \(t_1\) green and \(t_0\) red pears in
\(s_t\), \(t_0+t_1=t\),
\[
 \text{(9)}\quad P(s_t\mid h_\theta)=\prod_{i=1}^t P_{h_\theta}(q^{k_i}_i\mid s_{i-1})=\theta^{t_1}(1-\theta)^{t_0},
\]
which depends only on the two counts, and Bayes's theorem gives, up to a
normalizing constant,
\[
 \text{(10)}\quad P(h_\theta\mid s_t)=\mathrm{Norm}'\times\theta^{\lambda/2-1+t_1}(1-\theta)^{\lambda/2-1+t_0}
\]
\cite[pp.~633--634]{Romeijn2011InductiveLogicStatistics}. Here
\(P(h_\theta)\) and \(P(h_\theta\mid s_t)\) are densities over
\(\Theta\), as the integrals over \(d\theta\) show.

\paragraph{The controversy over priors.}
Romeijn locates most of the dispute over Bayesian method in the
determination and interpretation of the prior. Data are repeatable
events, so their probabilities can be read as frequencies or as some
other objective probability; a probability over hypotheses cannot be
read so and expresses epistemic uncertainty. Classical statistics
rejects that epistemic reading and keeps to probability as relative
frequency. Even granting the reading, nothing at the outset says which
hypothesis is right or a good candidate; objective criteria such as
maximum entropy or other forms of the principle of indifference have
been proposed, yet the starting point keeps a subjective element. The
strength of classical procedures is that they need no such prior
\cite[pp.~634--635]{Romeijn2011InductiveLogicStatistics}. These are
philosophical claims.

\section{Inductive logic with hypotheses}
\label{sec:p15-w02c-p15-029-rj-hypotheses}

The posterior of a Bayesian inference and the prediction of a Carnapian
logic are linked by the law of total probability,
\[
 \text{(11)}\quad P(q^1_{t+1}\mid s_t)=\int_0^1 P(q^1_{t+1}\mid h_\theta\cap s_t)P(h_\theta\mid s_t)\,d\theta .
\]
Under the multinomial hypotheses (7) the integrand's first factor is
\(\theta\), so
\[
 \text{(12)}\quad P(q^1_{t+1}\mid s_t)=\int_\Theta\theta P(h_\theta\mid s_t)\,d\theta=E[\theta]:
\]
the expected parameter value is a predictive probability. With the Beta
prior and posterior (10), the integral evaluates to
\[
 \text{(13)}\quad P(q^1_{t+1}\mid s_t)=\frac{t_1+\frac{\lambda}{2}}{t+\lambda},
\]
a Carnapian \(\lambda\)-rule (4) with \(n=2\)
\cite[p.~635]{Romeijn2011InductiveLogicStatistics}.

\begin{hbtheorem}{thm:p15-w02c-p15-029-rj-beta}{proved=annex:p15-w02c-p15-029-rj-beta}{Emma's farm and the \texorpdfstring{\(\lambda\)}{lambda}-rule}
Let \(0<\lambda<\infty\). On \(\Theta\times\{0,1\}^{\Omega}\) let the
joint law be: \(\theta\) has the density (8) with
\(\mathrm{Norm}=1/B(\tfrac\lambda2,\tfrac\lambda2)\), where
\(B(a,b)=\int_0^1\theta^{a-1}(1-\theta)^{b-1}d\theta\); and given
\(\theta\), the observations are independent with
\(P(q^1_i\mid h_\theta)=\theta\). Then for every \(t\ge0\) and every
\(s_t\) with counts \(t_0,t_1\):
\begin{enumerate}
\item (7) holds and the likelihood is (9);
\item \(P(s_t)=B(\tfrac\lambda2+t_1,\tfrac\lambda2+t_0)/B(\tfrac\lambda2,\tfrac\lambda2)>0\),
  and the posterior density is (10) with
  \(\mathrm{Norm}'=1/B(\tfrac\lambda2+t_1,\tfrac\lambda2+t_0)\);
\item (11) and (12) hold, and their common value is (13);
\item for \(\lambda=2\), (13) is \((t_1+1)/(t+2)\), the value \(c^*\) of
  \cref{thm:p15-w02b-p15-025-ml-cstar} with \(k=2\), \(w=1\),
  \(s=t\), \(s_M=t_1\).
\end{enumerate}
At \(\lambda=0\) the function (8) is not integrable on \([0,1]\), so
(8) defines no prior there, although (4) with \(\lambda=0\) is still a
formula.
\end{hbtheorem}

\begin{proofboundary}{def:p15-w02c-p15-029-rj-definetti-boundary}{De Finetti's representation as Romeijn uses it}
Romeijn states that, for any prior over the multinomial hypotheses, the
resulting predictive probabilities are exchangeable (2), and conversely
that any exchangeable predictive probability can be written as a
Bayesian inference over multinomial hypotheses with some prior, so that
de Finetti (1937) established a one-to-one correspondence between
exchangeability and Bayesian inference with multinomial hypotheses
\cite[pp.~635--636]{Romeijn2011InductiveLogicStatistics}. He gives no
proof. The handbook proves the representation in
\cref{thm:p15-binary-definetti} for the binary case and in
\cref{thm:p15-definetti-kernel} in general, with uniqueness of the
mixing law in \cref{thm:p15-directing-law-uniqueness}. Those theorems
concern infinite exchangeable sequences; the finite case differs, as
\cref{thm:p15-carnap-history-finite} shows by an exchangeable law on two
observations with no exchangeable extension.
\end{proofboundary}

\paragraph{Two readings of the representation.}
In the non-ampliative template, a Bayesian inference takes prior and
likelihoods as premises, and these leave a single probability
assignment over \(\mathcal H\times\mathcal Q\); a Carnapian logic is a
single assignment over \(\mathcal Q\). The representation shows that,
for predictions, the first can be reduced to the second
\cite[p.~636]{Romeijn2011InductiveLogicStatistics}. De Finetti welcomed
this: as a strict subjectivist who wanted to rid science of objective
chance, he used the theorem against multinomial and, more generally,
statistical hypotheses, since exchangeability is a predictive property
open to an empirical and subjective reading
\cite[p.~636]{Romeijn2011InductiveLogicStatistics}. Romeijn proposes
the inverse application. Since science persists in using statistical
hypotheses, the question is why; his answer, argued in his papers of
2004, 2005 and 2006, is that hypotheses help to pin down, indirectly,
assignments over \(\mathcal Q\) with the required predictive
properties. The representation then lets one capture relations between
observations by citing the hypotheses that may be true of them, a
natural extension of Carnapian logic
\cite[p.~636]{Romeijn2011InductiveLogicStatistics}. The two readings
use the same mathematics and differ in philosophy.

\paragraph{Bayesian inductive logic.}
Romeijn calls the resulting system \emph{Bayesian inductive logic}.
Its premises restrict the probability assignments over
\(\mathcal H\times\mathcal Q\); its conclusions are their probabilistic
consequences, derived by the axioms, often through Bayes's theorem; the
step is non-ampliative. For predictions alone the hypotheses can be
projected away, but the projected function may be very hard to define
by predictive properties alone. Bayesian statistics is captured
immediately: prior and likelihoods leave a unique model, and all
conclusions follow from the posterior
\cite[pp.~636--637]{Romeijn2011InductiveLogicStatistics}.

He names two precursors. Festa's presupposition view (1993) treats the
choice of \(\lambda\) as Romeijn treats the choice of hypotheses. And
Hintikka, replying to Kuipers in the Hintikka volume edited by Auxier
and Hahn (2006), holds that inductive inference rests on tacit
assumptions about the world and becomes in principle deductive once
they are spelled out. Romeijn identifies those tacit assumptions with
Carnap's symmetry principles and predictive properties, which a choice
of statistical hypotheses makes explicit
\cite[p.~637]{Romeijn2011InductiveLogicStatistics}.

\section{Neyman--Pearson testing}
\label{sec:p15-w02c-p15-029-rj-np}

A Neyman--Pearson test chooses between two fully specified probability
functions over the sample space. With \(\mathcal H=\{h_0,h_1\}\), each
\(h_j\) has a complete function \(P_{h_j}\) on \(\mathcal Q\), but
unlike Bayesian statistics the hypotheses are not in the probability
space: no probability is assigned to them and \(P(\cdot\mid h_j)\) is not
written \cite[pp.~637--638]{Romeijn2011InductiveLogicStatistics}.

\begin{hbdefinition}{def:p15-w02c-p15-029-rj-np-test}{Neyman--Pearson hypothesis test (Romeijn's Definition 2)}
A test function \(F\) on \(\mathcal Q\) is
\[
 \text{(14)}\quad F(s_t)=\begin{cases}1&\text{if }\dfrac{P_{h_1}(s_t)}{P_{h_0}(s_t)}>r,\\[1ex]0&\text{otherwise,}\end{cases}
\]
where \(P_{h_j}\) is the probability determined by \(h_j\). If \(F=1\),
the null hypothesis \(h_0\) is rejected; otherwise it is accepted for
the time being. In this simplified setting \(F\) is defined for each
set of sequences \(s_t\) separately, so each sample plan and sample size
needs its own test function. The test has
\begin{align*}
 \text{Significance}_F&=\alpha=\int_{\mathcal Q}F(s_t)P_{h_0}(s_t)\,ds_t,\\
 \text{Power}_F&=1-\beta=\int_{\mathcal Q}F(s_t)P_{h_1}(s_t)\,ds_t:
\end{align*}
\(\alpha\) is the type-I error of falsely rejecting \(h_0\), and
\(\beta\) the type-II error of falsely accepting it
\cite[p.~638]{Romeijn2011InductiveLogicStatistics}. Where
\(P_{h_0}(s_t)=0<P_{h_1}(s_t)\), the ratio in (14) is read here as
\(+\infty\); the source does not address this case.
\end{hbdefinition}

Romeijn reports that an optimal test minimizes significance and
maximizes power, and that Neyman and Pearson prove a decision optimal
in this sense for, and only for, likelihood-ratio test functions; he
refers to Barnett (1999) and to Neyman and Pearson's joint papers
\cite[p.~638]{Romeijn2011InductiveLogicStatistics}. The exact
statement, with its hypotheses, randomization at the threshold and the
converse, is \cref{thm:p15-neyman-pearson}.

\paragraph{Maria or Lisa.}
A pear of colour \(q^k\) comes from Maria's farm (\(h_0\)) or Lisa's
(\(h_1\)), whose colour compositions are
\[
 \begin{array}{c|ccc}
  & q^0 & q^1 & q^2\\\hline
  h_0 & 0.00 & 0.05 & 0.95\\
  h_1 & 0.40 & 0.30 & 0.30
 \end{array}
\]
The chosen test is
\[
 \text{(15)}\quad F(q^k)=\begin{cases}0&\text{if }k=2,\\1&\text{else,}\end{cases}
\]
with significance \(P_{h_0}(q^0\cup q^1)=0.05\) and power
\(P_{h_1}(q^0\cup q^1)=0.70\). A green pear gives \(F=1\), and Maria is
rejected as the grower with that significance and power
\cite[pp.~638--639]{Romeijn2011InductiveLogicStatistics}. In the
likelihood-ratio form (14) the source says this corresponds to a
threshold \(r\) between \(\tfrac{6}{19}\) and \(14\), for example
\(r=1\) \cite[p.~639]{Romeijn2011InductiveLogicStatistics}. With the
tabulated values the upper end is \(6\), not \(14\): for
\(6\le r<14\), (14) accepts \(h_0\) on a green pear.

\paragraph{Jeffreys's example.}
The decision depends on the whole distribution over the sample space,
not only on the probability of what was observed; optional stopping is
the familiar illustration. Romeijn uses instead an example he traces to
Jeffreys (1931, p.~357), a variant of which Hacking (1965) discusses;
footnote~4 credits the exact formulation to Jos Uffink
\cite[p.~639 and n.~4]{Romeijn2011InductiveLogicStatistics}. Replace
\(h_0\) by
\[
 \begin{array}{c|ccc}
  & q^0 & q^1 & q^2\\\hline
  h'_0 & 0.05 & 0.05 & 0.90\\
  h_1 & 0.40 & 0.30 & 0.30
 \end{array}
\]
Keeping significance \(P_{h'_0}(q^0)=0.05\) gives the test
\(F(q^k)=1\) if and only if \(k=0\), with power \(P_{h_1}(q^0)=0.40\). A
green pear \(q^1\) now leads to accepting \(h'_0\), although \(h_0\)
and \(h'_0\) give \(q^1\) the same probability. The null hypothesis
\(h'_0\) escapes rejection because it gives some probability to
\(q^0\), which did not occur, and so moves the rejection region
\cite[p.~640]{Romeijn2011InductiveLogicStatistics}. Romeijn quotes
Jeffreys's complaint that ``the null hypothesis can be rejected because
it fails to predict an event that never occurred''
\cite[p.~640]{Romeijn2011InductiveLogicStatistics}. The moral he draws
is that what a Neyman--Pearson procedure delivers depends on the
probability assignment over the entire sample space, not only on the
observation actually made. Seen from inductive logic the test may then
look like ``a remarkable procedure'', a second phrase he takes from
Jeffreys \cite[p.~640]{Romeijn2011InductiveLogicStatistics}. This is an
evaluation from the inferential standpoint; the reply to it is given
below.

\begin{hbtheorem}{thm:p15-w02c-p15-029-rj-np-tables}{proved=annex:p15-w02c-p15-029-rj-np-tables}{The two tests of the farm example}
Take one observation, the two tables above, and the convention of
\cref{def:p15-w02c-p15-029-rj-np-test} for a zero denominator.
\begin{enumerate}
\item For \(h_0\) against \(h_1\) the likelihood ratios of
  \(q^0,q^1,q^2\) are \(+\infty\), \(6\) and \(\tfrac{6}{19}\). The
  test (14) equals (15) exactly when \(\tfrac{6}{19}\le r<6\). Its
  significance is \(0.05\) and its power \(0.70\). Among the eight
  non-randomized tests (rejection regions
  \(R\subseteq\{q^0,q^1,q^2\}\)), those of significance at most
  \(0.05\) are \(\varnothing\), \(\{q^0\}\), \(\{q^1\}\),
  \(\{q^0,q^1\}\), with powers \(0\), \(0.40\), \(0.30\), \(0.70\); so
  (15) has the greatest power among them.
\item For \(h'_0\) against \(h_1\) the ratios are \(8\), \(6\) and
  \(\tfrac13\). The test rejecting exactly on \(q^0\) is (14) exactly
  when \(6\le r<8\); its significance is \(0.05\) and power \(0.40\). The
  non-randomized tests of significance at most \(0.05\) are
  \(\varnothing\), \(\{q^0\}\), \(\{q^1\}\), with powers \(0\),
  \(0.40\), \(0.30\); so it has the greatest power among them.
\item On the observation \(q^1\), the first test rejects its null and
  the second accepts its null, although
  \(P_{h_0}(q^1)=P_{h'_0}(q^1)=0.05\) and the alternative is the same.
\end{enumerate}
\end{hbtheorem}

\paragraph{Ampliative reading and Neyman and Pearson's own reading.}
For an ampliative logic it is natural to see an inferential step in the
test: \(F\) takes a sample and two functions \(P_{h_0},P_{h_1}\) to a
single function. The step is ampliative because both functions are
consistent with the data, so eliminating one is not deductive
\cite[p.~639]{Romeijn2011InductiveLogicStatistics}. Footnote~3 mentions
attempts to make such ampliative steps precise by default logic or by
preferential orderings of probability models, in particular Kyburg's
evidential probability (1974), discussed more recently by Wheeler
(2006); Romeijn sets these aside because they concern inferences among
hypotheses and from hypotheses to predictions, not from data
\cite[p.~639 n.~3]{Romeijn2011InductiveLogicStatistics}. His answer
to the evaluation after Jeffreys's example, which follows it
immediately in the source, is that Neyman--Pearson statistics was never
meant as a disguised inference: it decides between hypotheses and attaches error rates to the
decision; Neyman and Pearson explicitly rejected an inferential reading,
and significance and power are probabilities of data given hypotheses,
not of hypotheses given data \cite[p.~640]{Romeijn2011InductiveLogicStatistics}.

\section{The Neyman--Pearson test as an inference}
\label{sec:p15-w02c-p15-029-rj-np-inference}

Romeijn next asks whether the test can be turned into an inference of
Bayesian inductive logic, calling the exercise, with self-irony,
statistical yoga. He gives two reasons for doing it anyway: it shows how
statistics relates to inductive logic, and practitioners are tempted
to read the significance as the probability that the null hypothesis
is true, a reading he calls emphatically wrong but intuitively
appealing, which the construction shows to treat the test as a
non-ampliative entailment \cite[p.~640]{Romeijn2011InductiveLogicStatistics}.

\begin{hbdefinition}{def:p15-w02c-p15-029-rj-np-model}{Interval priors and coarse data}
Build the space \(\mathcal H\times\mathcal Q\) with the functions
\(P_{h_j}\) on the sample spaces \(\langle h_j,\mathcal Q\rangle\). For
the prior take \(P(h_0)\in(l,u)\), that is \(l<P(h_0)<u\), with
\(P(h_0)+P(h_1)=1\); write \(\underline P(h_j)=\min P(h_j)\) and
\(\overline P(h_j)=\max P(h_j)\) over the set. Coarse-grain each
observation to the value of the test function,
\[
 f^j=\{s_t:\ F(s_t)=j\},
\]
so that the error rates become likelihoods,
\(P(f^1\mid h_0)=\alpha\) and \(P(f^0\mid h_1)=\beta\), and Bayes's
theorem gives the set of posteriors through
\[
 \frac{P(h_1\mid f^j)}{P(h_0\mid f^j)}=\frac{P(f^j\mid h_1)P(h_1)}{P(f^j\mid h_0)P(h_0)}
\]
\cite[pp.~640--641]{Romeijn2011InductiveLogicStatistics}. Over an open
interval of priors the bounds are not attained; \(\min\) and \(\max\)
are read below as infimum and supremum.
\end{hbdefinition}

Romeijn's claims about this model are these. The quality of the test is
reflected in the posteriors: after \(f^1\), ``generally''
\(\underline P(h_0\mid f^1)<\underline P(h_0)\) and
\(\overline P(h_0\mid f^1)<\overline P(h_0)\). The prior interval
\((0,1)\) yields the posterior interval \((0,1)\), mirroring classical
statistics' refusal of a prior; but if
\[
 l\ge\frac{\beta}{\beta+1-\alpha},\qquad u\le\frac{1-\beta}{1-\beta+\alpha},
\]
then every prior in \((l,u)\) gives
\(\overline P(h_0\mid f^1)<\tfrac12<\underline P(h_1\mid f^1)\) and
\(\underline P(h_0\mid f^0)>\tfrac12>\overline P(h_1\mid f^0)\), so
\(F=1\) tilts every function in the set toward \(h_1\) and \(F=0\)
toward \(h_0\) \cite[p.~641]{Romeijn2011InductiveLogicStatistics}. In
the farm example, \(f^1=q^0\cup q^1\), \(f^0=q^2\), \(\alpha=0.05\),
\(\beta=0.30\), and the prior interval \((0.24,0.93)\) yields
\(P(h_0\mid f^0)\in(0.50,0.98)\) and \(P(h_0\mid f^1)\in(0.02,0.50)\)
\cite[p.~641]{Romeijn2011InductiveLogicStatistics}. The source's display
of the second likelihood reads
\(P(f^0\mid h_1)=P(q^0\cup q^1\mid h_1)=\beta=0.30\); since \(f^0=q^2\),
the correct middle term is \(P(q^2\mid h_1)\), which is \(0.30\), while
\(P(q^0\cup q^1\mid h_1)=0.70\).

\begin{hbtheorem}{thm:p15-w02c-p15-029-rj-interval-prior}{proved=annex:p15-w02c-p15-029-rj-interval-prior}{Exact form of the interval-prior rendering}
Let \(0<\alpha<1\) and \(0<\beta<1\), and for a prior \(p=P(h_0)\in[0,1]\) put
\begin{align*}
 g_1(p)&=P(h_0\mid f^1)=\frac{\alpha p}{\alpha p+(1-\beta)(1-p)},\\
 g_0(p)&=P(h_0\mid f^0)=\frac{(1-\alpha)p}{(1-\alpha)p+\beta(1-p)}.
\end{align*}
\begin{enumerate}
\item \(g_0\) and \(g_1\) are continuous and strictly increasing on
  \([0,1]\), with value \(0\) at \(0\) and \(1\) at \(1\). For
  \(0\le l<u\le1\) the set of posteriors from priors in \((l,u)\) is the
  open interval \((g_j(l),g_j(u))\); in particular \((0,1)\) goes to
  \((0,1)\).
\item \(g_1(p)<\tfrac12\) if and only if \(p<\frac{1-\beta}{1-\beta+\alpha}\),
  and \(g_0(p)>\tfrac12\) if and only if \(p>\frac{\beta}{\beta+1-\alpha}\).
  Hence, for a nonempty interval \((l,u)\), every prior in it satisfies
  both \(P(h_0\mid f^1)<\tfrac12<P(h_1\mid f^1)\) and
  \(P(h_0\mid f^0)>\tfrac12>P(h_1\mid f^0)\) if and only if
  \(l\ge\frac{\beta}{\beta+1-\alpha}\) and
  \(u\le\frac{1-\beta}{1-\beta+\alpha}\). Such an interval exists if and
  only if \(\alpha+\beta<1\).
\item \(g_1(p)<p<g_0(p)\) for every \(p\in(0,1)\) if and only if
  \(\alpha+\beta<1\), that is, the power exceeds the significance. Under
  this condition \(\inf g_1=g_1(l)<l\) and \(\sup g_1=g_1(u)<u\) over
  \((l,u)\) with \(0<l<u<1\); if \(\alpha+\beta>1\), both inequalities
  reverse.
\item For \(\alpha=0.05\), \(\beta=0.30\): the two bounds are
  \(\frac{6}{25}=0.24\) and \(\frac{14}{15}\approx0.933\). For the prior
  interval \((\frac{6}{25},\frac{14}{15})\) the posterior intervals are
  \(P(h_0\mid f^0)\in(\tfrac12,\tfrac{133}{136})\) and
  \(P(h_0\mid f^1)\in(\tfrac{3}{136},\tfrac12)\), with
  \(\tfrac{133}{136}\approx0.978\) and \(\tfrac{3}{136}\approx0.022\).
  With the upper end rounded to \(0.93\), the upper posterior ends are
  about \(0.977\) and \(0.487\).
\end{enumerate}
\end{hbtheorem}

So the source's upper value \(0.50\) for \(P(h_0\mid f^1)\) belongs to
the exact bound \(\tfrac{14}{15}\), of which \(0.93\) is a rounding; and
its ``generally'' is the condition that the power exceed the
significance.

\paragraph{How far the imitation goes.}
Depending on the prior, one may say that the Bayesian inference
replicates the test: if probabilities over hypotheses express
preference, \(f^0\) favours \(h_0\) and \(f^1\) favours \(h_1\). The
inference fits the entailment relation of
\cref{def:p15-w02c-p15-029-rj-templates}: the premises are the set of
priors over \(\mathcal H\) coupled with the full assignments over
\(\langle h_j,\mathcal Q\rangle\), the conclusions the set of
posteriors, and because the latter are computed from the former by the
axioms, every premise function lies in the conclusion set
\cite[p.~642]{Romeijn2011InductiveLogicStatistics}. The imitation is
far from perfect. A Bayesian inference ends in a probability function,
while a test ends in a binary decision; rules such as choosing \(h_0\)
when \(P(h_0\mid s_t)>\tfrac12\) can be added, as in Levi (1980) or the
literature on rational acceptance such as Douven (2002), but that
bivalence then comes from a decision-theoretic add-on, not from the
inference. And advocates of classical statistics may reject any prior,
however unspecific, on the ground that probabilities are long-run
frequencies, which hypotheses lack \cite[p.~642]{Romeijn2011InductiveLogicStatistics}.

\paragraph{Likelihood principle or total evidence.}
Following Hacking (1965) and Edwards (1972), Jeffreys's example shows
that Neyman--Pearson tests violate the likelihood principle, since they
depend on the whole sample-space distribution. In the Bayesian
rendering the likelihood principle is respected, but the conditioning is
on \(f^j\), not on \(q^k\). Romeijn therefore offers a different
diagnosis: the procedure takes the observations through a coarse
partition, so it can as well be said to violate the principle of total
evidence \cite[pp.~642--643]{Romeijn2011InductiveLogicStatistics}. The
next theorem makes the two diagnoses exact for the farm example.

\begin{hbtheorem}{thm:p15-w02c-p15-029-rj-total-evidence}{proved=annex:p15-w02c-p15-029-rj-total-evidence}{Jeffreys's example under fine and coarse conditioning}
Let the prior give the null hypothesis probability \(p\in(0,1)\) and
\(h_1\) probability \(1-p\), in each of the two problems of
\cref{thm:p15-w02c-p15-029-rj-np-tables}.
\begin{enumerate}
\item Conditioning on the full observation \(q^1\) gives the same
  posterior in both problems:
  \(P(h_0\mid q^1)=P(h'_0\mid q^1)=\frac{0.05p}{0.05p+0.30(1-p)}\).
\item Conditioning on the coarse observation containing \(q^1\) gives
  \[
   P(h_0\mid f^1)=\frac{0.05p}{0.05p+0.70(1-p)},
  \]
  \[
   P(h'_0\mid f^0)=\frac{0.95p}{0.95p+0.60(1-p)}
  \]
  in the first and the second problem, where \(f^0=q^1\cup q^2\) in the
  second. These differ for every \(p\in(0,1)\);
  the first is below \(p\) and the second above it. At \(p=\tfrac12\) the
  three values are \(\tfrac17\), \(\tfrac1{15}\) and \(\tfrac{19}{31}\).
\end{enumerate}
\end{hbtheorem}

\section{Fisher's parameter estimation}
\label{sec:p15-w02c-p15-029-rj-fisher}

The second classical procedure is maximum likelihood estimation, which
Romeijn attributes to Fisher with reference to the 1956 book and to
Barnett (1999). The estimator picks the best of a larger, possibly
infinite, set of hypotheses, using the probabilities they give to
points of the sample space \cite[p.~643]{Romeijn2011InductiveLogicStatistics}.

\begin{hbdefinition}{def:p15-w02c-p15-029-rj-mle}{Maximum likelihood estimation and the confidence region (Romeijn's Definition 3)}
Let \(\mathcal H=\{h_\theta:\theta\in\Theta\}\) and let \(\mathcal Q\)
be the sample space. The maximum likelihood estimator is the function
on the sample space
\[
 \text{(16)}\quad \hat\theta(s_t)=\{\theta\in\Theta:\ \forall h_{\theta'}\,\bigl(P_{h_{\theta'}}(s_t)\le P_{h_\theta}(s_t)\bigr)\},
\]
a set of parameter values, typically a singleton; the associated best
hypothesis is \(h_{\hat\theta(s_t)}\), or \(h_{\hat\theta}\). With
\(\Theta=[0,1]\) and true value \(\theta\), a region of sample space in
which the estimator is not far off covers \(1-\epsilon\) of
\(P_{h_\theta}\):
\[
 \text{(17)}\quad \mathrm{Conf}_{1-\epsilon}(\theta)=\Bigl\{\hat\theta:\ |\hat\theta-\theta|<\Delta\ \text{and}\ \int_{\theta-\Delta}^{\theta+\Delta}P_{h_\theta}(\hat\theta)\,d\hat\theta=1-\epsilon\Bigr\}
\]
\cite[p.~643]{Romeijn2011InductiveLogicStatistics}.
\end{hbdefinition}

The interval \(\hat\theta\in[\theta-\Delta,\theta+\Delta]\) has a
frequentist reading: in a series of estimations the fraction falling
further than \(\Delta\) from the mark tends to \(\epsilon\), and a
smaller region means a more reliable estimate. Its drawback, Romeijn
notes, is that it is defined through the unknown true value
\cite[p.~643]{Romeijn2011InductiveLogicStatistics}. For the multinomial
hypotheses (7) the likelihood is \(\theta^{t_1}(1-\theta)^{t_0}\),
maximal at
\[
 \text{(18)}\quad \hat\theta(s_t)=\frac{t_1}{t};
\]
the probability that \(t_1/t\in[\theta-\Delta,\theta+\Delta]\) is said
to increase with the length of the data by the law of large numbers, so
that at fixed probability \(1-\epsilon\) the interval shrinks
\cite[pp.~643--644]{Romeijn2011InductiveLogicStatistics}. The maximizer
for \(t\ge1\), its unbiasedness, variance and the Chebyshev bound
\(\theta(1-\theta)/(t\Delta^2)\) on the miss probability are
\cref{thm:p15-fe-binomial}. Read literally, ``increases'' is too
strong.

\begin{hbtheorem}{thm:p15-w02c-p15-029-rj-mle-coverage}{proved=annex:p15-w02c-p15-029-rj-mle-coverage}{Coverage tends to one, not monotonically}
Under \(h_\theta\) of (7), let
\(c_t=P_{h_\theta}\bigl(t_1/t\in[\theta-\Delta,\theta+\Delta]\bigr)\).
\begin{enumerate}
\item If \(t=0\), every \(\theta\in\Theta\) belongs to \(\hat\theta(s_0)\).
\item For every \(\Delta>0\), \(c_t\ge1-\theta(1-\theta)/(t\Delta^2)\),
  so \(c_t\to1\).
\item For \(\theta=\tfrac12\) and \(\Delta=\tfrac14\),
  \(c_4=\tfrac78\) and \(c_5=\tfrac58\); so \(c_t\) is not monotone in
  \(t\).
\end{enumerate}
\end{hbtheorem}

Romeijn reads estimation as the procedural analogue of an ampliative
step: from the set \(\mathcal H\), the data select one \(h_{\hat\theta}\)
as conclusion. The next two sections look for non-ampliative
representations \cite[p.~644]{Romeijn2011InductiveLogicStatistics}.

\section{Estimation in inductive logic}
\label{sec:p15-w02c-p15-029-rj-estimation}

Romeijn names two routes to a non-ampliative estimation. One, fiducial
inference, gives a probability over hypotheses without a prior and is
treated in \cref{sec:p15-w02c-p15-029-rj-fiducial}. The other relates
estimation to the logics already introduced
\cite[p.~644]{Romeijn2011InductiveLogicStatistics}.

\paragraph{Estimation as a Carnapian predictive system.}
An estimate can be used to predict. After \(s^{000101}\), one chooses
\(h_{1/3}\) and predicts green with probability \(\tfrac13\). The
estimator then acts as a predictive system,
\(P(q^k_{t+1}\mid s_t)=P_{\hat\theta(s_t)}(q^k_{t+1})\), where
\(P_{\hat\theta(s_t)}\) is the function of \(h_{\hat\theta(s_t)}\); so
Fisher's estimator is captured by a single function over the sample
space from which estimates follow non-ampliatively. For red and green
pears, the \(\lambda\)-rule (4) with \(\lambda=0\) predicts the observed
relative frequencies, which by (18) are the maximum likelihood
estimates, so on each set of possible observations
\(\{s^{k_1\ldots k_t}:k_i=0,1\}\) the rule predicts as Fisher estimates
\cite[pp.~644--645]{Romeijn2011InductiveLogicStatistics}. Footnote~5
records the cost. After \(s^{000}\) the probability of green next is
\(0\), so \(P(s^{0001})=0\) and the standard axioms leave
\(P(q^0_5\mid s^{0001})\) undefined, whereas following Fisher requires
\(\tfrac34\). The axiomatization must then change, for instance to one
with primitive conditional probability as in R\'enyi (1970) and a
chapter of the volume that the footnote numbers 15; alternatively one
may restrict estimation to data showing more than one property
\cite[p.~645 n.~5]{Romeijn2011InductiveLogicStatistics}.

\begin{hbtheorem}{thm:p15-w02c-p15-029-rj-lambda-zero}{proved=annex:p15-w02c-p15-029-rj-lambda-zero}{The \texorpdfstring{\(\lambda=0\)}{lambda = 0} rule and Fisher's estimate}
Let \(n=2\).
\begin{enumerate}
\item For \(t\ge1\), (4) with \(\lambda=0\) gives
  \(P(q^1_{t+1}\mid s_t)=t_1/t=\hat\theta(s_t)\), the unique element of
  (16) under (7); for \(s^{000101}\) this is \(\tfrac13\).
\item Let \(P\) satisfy (4) with \(\lambda=0\) for \(t\ge1\) whenever
  the conditioning event has positive probability, and let
  \(P(q^0_1)>0\). Then \(P(s^{000})=P(q^0_1)>0\),
  \(P(q^1_4\mid s^{000})=0\) and \(P(s^{0001})=0\); so the ratio
  definition leaves \(P(q^0_5\mid s^{0001})\) undefined, while
  \(\hat\theta(s^{0001})=\tfrac14\) gives red the value \(\tfrac34\).
\end{enumerate}
\end{hbtheorem}

Romeijn does not pursue this route. Already for multinomial hypotheses
the corresponding function over the sample space is not obvious; for
more complex hypotheses and estimators it may be less natural still,
and principles or predictive properties motivating it will be hard to
find \cite[p.~645]{Romeijn2011InductiveLogicStatistics}.

\paragraph{Estimation in Bayesian inductive logic.}
Bayesian inductive logic suits estimation better: both procedures
consider a set of hypotheses, seek the best-fitting one and are guided
by likelihoods. From the posterior one obtains the expectation
\(E[\theta]=\int_\Theta\theta P(h_\theta\mid s_t)\,d\theta\), as in (12),
a function from data to a preferred parameter value which depends on the
prior but is analogous to the maximum likelihood estimator
\cite[p.~645]{Romeijn2011InductiveLogicStatistics}.

\begin{hbdefinition}{def:p15-w02c-p15-029-rj-credal}{Credal interval}
In analogy to (17), Romeijn defines from the posterior
\[
 \mathrm{Cred}_{1-\epsilon}(s_t)=\Bigl\{\theta:\ |\theta-E[\theta]|<\Delta\ \text{and}\ \int_{E[\theta]-\Delta}^{E[\theta]+\Delta}P(h_\theta\mid s_t)\,d\theta=1-\epsilon\Bigr\},
\]
so that
\[
 \text{(19)}\quad P\bigl(\{\theta:\theta\in\mathrm{Cred}_{1-\epsilon}(s_t)\}\mid s_t\bigr)=1-\epsilon :
\]
the posterior probabilities of the corresponding \(h_\theta\) add up to
\(1-\epsilon\) \cite[pp.~645--646]{Romeijn2011InductiveLogicStatistics}.
\end{hbdefinition}

Romeijn suggests that (19) improves on (17): the confidence interval
says how far an estimate is off on average but licenses no inference
about how far the particular estimate lies from the true value, whereas
the credal interval does allow such an inferential reading
\cite[p.~646]{Romeijn2011InductiveLogicStatistics}. The two kinds of
interval answer different questions;
\cref{thm:p15-confidence-credibility-boundary} computes, for a binomial
sample, a posterior content of an exact confidence interval that differs
from its confidence level. Two differences remain. Classical estimation
is all-or-nothing (accept, reject, one best estimate), while the
Bayesian result is a posterior over all hypotheses, and this remains
problematic as for testing. And the chosen prior fixes every
expectation and every credal interval; that dependence is avoided only by assuming
enough data to swamp the prior or, roughly equivalently, a prior smooth
relative to the likelihoods \cite[p.~646]{Romeijn2011InductiveLogicStatistics}.

\section{Fiducial probability}
\label{sec:p15-w02c-p15-029-rj-fiducial}

The problem of choosing a prior led Fisher (1930, 1935, 1956) to the
fiducial argument, which yields a probability over hypotheses without a
prior. Romeijn calls it controversial and of limited applicability,
refers to Hacking (1965), Seidenfeld (1979) and Barnett (1999), and
offers only a sketch \cite[p.~646]{Romeijn2011InductiveLogicStatistics}.

\paragraph{The fiducial interval.}
Romeijn introduces the argument through the confidence intervals (17) of
\cref{sec:p15-w02c-p15-029-rj-fisher}: in some cases the data also
yield a region of parameter values expected to contain the true value
\(\theta\). In general terms, one fixes a set \(R\) of parameter values
under which the observed data are not too improbable, and then asserts
that the true parameter value most probably belongs to \(R\)
\cite[p.~646]{Romeijn2011InductiveLogicStatistics}. The interval (20)
is his specific form of \(R\). Swapping the roles of \(\theta\) and \(\hat\theta\) in (17) gives
\[
 \text{(20)}\quad \mathrm{Fid}_{1-\epsilon}(s_t)=\Bigl\{\theta:\ |\theta-\hat\theta(s_t)|<\Delta\ \text{and}\ \int_{\hat\theta-\Delta}^{\hat\theta+\Delta}P_{h_\theta}(\hat\theta)\,d\theta=1-\epsilon\Bigr\},
\]
where the integral now runs over parameter values, not data. Each
sample receives a fiducial interval of parameter values considered good
candidates \cite[pp.~646--647]{Romeijn2011InductiveLogicStatistics}.
The integral is a probability only if the roles can be swapped, which
needs
\[
 P_{h_\theta}(\hat\theta+\delta)=P_{h_{\theta-\delta}}(\hat\theta)\quad\text{for all }\delta,
\]
so that the distribution of the estimator given the parameter can be
read as a distribution of the parameter given the estimator. Then the
fiducial interval reads as a probability like (19),
\[
 \text{(21)}\quad P\bigl(\{\theta:\theta\in\mathrm{Fid}_{1-\epsilon}(s_t)\}\mid s_t\bigr)=1-\epsilon;
\]
otherwise it cannot, at least without further consideration, be read as
the probability that the true value lies near the estimate
\cite[p.~647]{Romeijn2011InductiveLogicStatistics}. The swap condition
says that \(P_{h_\theta}(x)\) depends on \(x\) and \(\theta\) only
through \(x-\theta\): putting \(\delta=x-\hat\theta\) and then
\(\hat\theta=x-\theta\) shows \(P_{h_\theta}(x)=P_{h_0}(x-\theta)\) when
\(0\in\Theta\), the defining property of a location family.

\paragraph{Functional models.}
Romeijn calls the requirement strong: two distinct functions, both
written \(P_{h_\theta}(\hat\theta)\), one of \(\theta\) and one of
\(\hat\theta\), must coincide. Dawid and Stone (1982) give a far more
general formulation: the fiducial argument needs a functional model of
the problem that is smoothly invertible
\cite[p.~647]{Romeijn2011InductiveLogicStatistics}.

\begin{hbdefinition}{def:p15-w02c-p15-029-rj-functional-model}{Pivotal quantity and functional model}
A \emph{pivotal quantity} is an estimator function \(\hat\theta(s_t)\)
related to the parameter \(\theta\) and an error term \(\omega\) by
\[
 \hat\theta(s_t)=f(\theta,\omega),
\]
with \(\theta\) the systematic and \(\omega\) the stochastic component of
the process producing the data. A probability \(P(\omega)\) over error
terms is assumed, so that
\[
 \text{(22)}\quad P(\hat\theta(s_t)\mid h_\theta)=P(\{\omega:\ f(\theta,\omega)=\hat\theta\}).
\]
The function \(f\) is assumed invertible, with \(f^{-1}(\hat\theta,\omega)=\theta\),
and error terms and hypotheses are assumed independent,
\[
 \text{(23)}\quad P(h_\theta,\omega)=P(h_\theta)P(\omega),
\]
so every parameter value carries the same distribution of error terms.
The joint law is then a Bayesian network with arrows from \(h_\theta\)
and from \(\omega\) into \(\hat\theta\) (Romeijn refers to Neapolitan
2003). Observing \(s_t\) fixes \(\hat\theta(s_t)\); conditioning on it,
the network structure and the determinism of \(f\) make \(\omega\) and
\(h_\theta\) perfectly correlated, each \(\omega\) going with the unique
\(\theta=f^{-1}(\hat\theta,\omega)\); and since, Romeijn adds, the
observation does not itself affect the probability of \(\omega\),
\[
 \text{(24)}\quad P(h_\theta\mid\hat\theta(s_t))=P(\{\omega:\ f^{-1}(\hat\theta,\omega)=\theta\}),
\]
the inverse of (22): the distribution over \(\omega\) is transferred onto
\(h_\theta\) by \(f^{-1}\) \cite[pp.~647--648]{Romeijn2011InductiveLogicStatistics}.
\end{hbdefinition}

Romeijn finds the result surprising: no prior is assumed, yet something
like a posterior is derived, and the inference is non-ampliative. In
his description, \(P(h_\theta)\) may be any convex combination of
members of a set of functions over \(\theta\), while \(P(h_\theta\mid s_t)\)
is a definite member of that set; he regrets that the argument runs only
under such strict conditions \cite[p.~648]{Romeijn2011InductiveLogicStatistics}.
The premise that observing \(s_t\) leaves the distribution of
\(\omega\) unchanged carries the weight of (24); it does not follow from
(23), and it constrains the prior. The next theorem makes this exact in
a finite setting.

\begin{hbtheorem}{thm:p15-w02c-p15-029-rj-fiducial-transfer}{proved=annex:p15-w02c-p15-029-rj-fiducial-transfer}{When the fiducial transfer (24) holds}
Let \(\Theta\) and the set \(W\) of error terms be finite, let \(P\) be
a probability on \(\Theta\times W\) satisfying (23), and let
\(\hat\theta=f(\theta,\omega)\) where, for each value \(x\) and each
\(\omega\), exactly one \(\theta\) has \(f(\theta,\omega)=x\), written
\(f^{-1}(x,\omega)\).
\begin{enumerate}
\item (22) holds for every \(\theta\) with \(P(h_\theta)>0\).
\item If \(P(\hat\theta=x)>0\) and \(P(\omega\mid\hat\theta=x)=P(\omega)\)
  for every \(\omega\), then (24) holds at \(x\) for every \(\theta\).
\item (23) and invertibility do not imply (24). For
  \(\Theta=W=\{0,1\}\), \(f(\theta,\omega)=\theta+\omega\bmod2\),
  \(P(\omega=0)=P(\omega=1)=\tfrac12\) and \(P(h_0)=1\):
  \(P(h_0\mid\hat\theta=x)=1\), while the right side of (24) is
  \(\tfrac12\), for both values \(x\).
\item Let \(\Theta=W=\mathbb Z_m\), \(f(\theta,\omega)=\theta+\omega\bmod m\).
  If the prior is uniform, then for every error law the premise of
  item~2 holds, and (24) gives \(P(h_\theta\mid\hat\theta=x)=P(\omega=x-\theta)\),
  which is also the Bayesian posterior. If the error law is uniform, the
  premise of item~2 holds for every \(x\) if and only if the prior is
  uniform.
\end{enumerate}
\end{hbtheorem}

In the cyclic model, therefore, the fiducial posterior is not obtained
without a prior: the premise of (24) holds exactly for the uniform
prior when the error law is uniform. Romeijn's ``any convex
combination'' describes the premise set before that premise is added.

\section{Conclusion, chronology and interpretations}
\label{sec:p15-w02c-p15-029-rj-assessment}

Romeijn concludes that he has set up a frame in which inductive logic
and statistics may be unified; that the discussion of Bayesian
inference extended the language of inductive logic with statistical
hypotheses; and that the two classical procedures, though best
understood as ampliative, can be modelled at least partly in the
extended logic. He hopes that the relation between Carnapian logic and
Bayesian statistics will stimulate research, names earlier work by
Skyrms (1991, 1993, 1996) and Festa (1993), and recalls that his
dissertation (2005) argued for the close relation of an inductive logic
with statistical hypotheses to Bayesian statistical inference
\cite[pp.~648--649]{Romeijn2011InductiveLogicStatistics}.

\paragraph{Chronology of the works the chapter relies on.}
The chapter is systematic, but the dates it cites fix an order:
Jeffreys's \emph{Scientific Inference} (1931), the source of the
rejected-null example; Fisher on inverse probability (1930) and the
fiducial argument (1935); Johnson (1932) and de Finetti (1937) on
exchangeability and prediction; Carnap's \emph{Logical Foundations of
Probability} (1950) and \emph{Continuum of Inductive Methods} (1952);
Fisher's \emph{Statistical Methods and Scientific Inference} (1956);
Hacking (1965); Neyman and Pearson's \emph{Joint Statistical Papers}
(1967); R\'enyi (1970); Edwards (1972); Kyburg (1974); Seidenfeld (1979);
Levi (1980); Dawid and Stone (1982) and Zabell (1982); Skyrms (1991,
1993, 1996); Festa (1993, 1996); Hailperin (1996); Barnett (1999);
Maher (2000); Douven (2002); Howson (2003), Press (2003) and Neapolitan
(2003); Romeijn (2004, 2005, 2006); Auxier and Hahn (2006) and Wheeler
(2006); and Bandyopadhyay and Forster (2009) and Haenni and others
(2009) \cite[pp.~649--650]{Romeijn2011InductiveLogicStatistics}. The
Paris--Waterhouse manuscript of 2008 is listed as unpublished.

\paragraph{Interpretations that the chapter keeps apart.}
\begin{enumerate}
\item \emph{Logical and subjectivist readings of exchangeability.} For
  Carnap the predictive properties follow from logical probability and
  so bear on Hume's problem; for de Finetti they had no objective ground,
  and the representation theorem served to eliminate chances. Romeijn
  inverts that use, taking hypotheses as aids to fixing predictions
  (\cref{sec:p15-w02c-p15-029-rj-carnap,sec:p15-w02c-p15-029-rj-hypotheses}).
\item \emph{Tests as decisions or as inferences.} Neyman and Pearson
  read the test as a decision with error rates; Romeijn constructs an
  inferential imitation and keeps the difference of binary decision and
  posterior (\cref{sec:p15-w02c-p15-029-rj-np-inference}).
\item \emph{Two diagnoses of Jeffreys's example.} Hacking and Edwards:
  violation of the likelihood principle. Romeijn: conditioning on a
  coarse partition, a violation of total evidence
  (\cref{thm:p15-w02c-p15-029-rj-total-evidence}).
\item \emph{Confidence, credal and fiducial intervals.} The confidence
  region is a sampling statement indexed by the unknown value; the
  credal interval is a posterior statement relative to a prior; the
  fiducial interval is a posterior-like statement whose probabilistic
  reading depends on the swap condition or an invertible functional
  model.
\end{enumerate}

\paragraph{Limitations.}
The chapter states several results without proof: the derivation of
(4) from (1)--(3), de Finetti's representation, and the Neyman--Pearson
optimality; their exact forms are cited above. Its sample space
excludes continua of individuals and of properties, yet Sections~9--11
integrate over a continuous parameter and a continuous estimator, so
\(P_{h_\theta}(\hat\theta)\) in (17) and (20) must be read as a density
of the estimator's distribution, which the chapter does not construct.
The printed table bound \(14\) for \(r\), the middle term of the
\(\beta\) display and the rounding \(0.93\) are corrected above. And the
fiducial transfer (24) rests on a premise the chapter states in words
but does not separate from (23).

\section{Romeijn's account beside Carnap's Section 98}
\label{sec:p15-w02c-p15-029-rj-comparison}

\Cref{sec:p15-w02b-p15-025-carnap98} sets Carnap's 1950 account of
estimation beside Fisher, Neyman--Pearson and Wald. Romeijn's chapter
treats the same pair of disciplines sixty years later, and the two
treatments can be compared point by point. This comparison concerns
induction and statistical inference; Romeijn's chapter contains no
measure of confirmation or evidential relevance of the kind compared in
\cref{sec:p15-w02b-p15-025-sources}, so it has no counterpart there.
\begin{enumerate}
\item \emph{Where hypotheses live.} Carnap's estimate is a mean weighted
  by degrees of confirmation of statements about the parameter, within a
  language of individuals and predicates
  (\cref{def:p15-carnap-c-mean}). Romeijn adds the hypotheses to the
  probability space, \(\mathcal H\times\mathcal Q\), and reads Carnap's
  restriction to \(\mathcal Q\) through de Finetti's projection.
\item \emph{Maximum likelihood.} Carnap measures the distance between
  the \(c^*\)-value and the observed frequency
  (\cref{thm:p15-w02b-p15-025-ml-cstar}). Romeijn identifies the
  observed frequency with the \(\lambda=0\) rule
  (\cref{thm:p15-w02c-p15-029-rj-lambda-zero}) and the binary \(c^*\)
  value with \(\lambda=2\) and the uniform prior
  (\cref{thm:p15-w02c-p15-029-rj-beta}, item~4); so the two compared
  values are two members of one \(\lambda\)-family.
\item \emph{Confidence.} Carnap treats the confidence coefficient as a
  direct inference relative to a hypothesis about the parameter, not a
  confirmation relative to the sample. Romeijn says the same of (17),
  that it is defined through the unknown true value, and offers the
  credal and fiducial intervals as statements relative to the sample.
\item \emph{Tests and priors.} Carnap's Section~98, as presented in
  \cref{sec:p15-w02b-p15-025-carnap98}, discusses Neyman's confidence
  method and Wald's minimax rule and reads the motive of the classical
  methods as rejection of the principle of indifference. Romeijn treats
  Neyman--Pearson testing directly, credits classical procedures with
  needing no subjective prior, and shows that an inferential imitation
  needs an interval prior with the bounds of
  \cref{thm:p15-w02c-p15-029-rj-interval-prior}.
\item \emph{Normative stance.} Carnap hoped one confirmation function
  would give a single method for all estimation problems. Romeijn
  disclaims norms for statistical inference and claims only that the
  logical reading clarifies the procedures.
\end{enumerate}

\section{Sources for this chapter}
\label{sec:p15-w02c-p15-029-rj-sources}

Romeijn, ``Inductive Logic and Statistics'', \emph{Handbook of the
History of Logic}, vol.~10, \emph{Inductive Logic}, pp.~625--650
\cite{Romeijn2011InductiveLogicStatistics}: sec.~1, pp.~625--626;
sec.~2, pp.~626--628 with n.~1; sec.~3, pp.~628--630; sec.~4,
pp.~630--632 with n.~2, equations (1)--(5); sec.~5, pp.~632--635,
Definition~1 and (6)--(10); sec.~6, pp.~635--637, (11)--(13); sec.~7,
pp.~637--640, Definition~2, (14)--(15), nn.~3--4; sec.~8, pp.~640--643;
sec.~9, pp.~643--644, Definition~3, (16)--(18); sec.~10, pp.~644--646,
(19), n.~5; sec.~11, pp.~646--648, (20)--(24); sec.~12, pp.~648--649;
bibliography, pp.~649--650. The works cited by Romeijn are reported as
he cites them and are not presented here. The theorems of this chapter
are not in the source; they are proved in
\cref{sec:annex-p15-w02c-p15-029-romeijn}.

\index{inductive logic!and statistics}
\index{statistical hypothesis!in inductive logic}
\index{Bayesian inductive logic}
\index{lambda-rule@\(\lambda\)-rule!instantial relevance}
\index{Neyman--Pearson test!as inference}
\index{fiducial argument!functional model}
\index{credal interval}
\index{total evidence!and Neyman--Pearson tests}
